% concatenated from main.tex
\documentclass[1p,final]{elsarticle}


\usepackage{lipsum}


\makeatletter
\def\ps@pprintTitle{%
	\let\@oddhead\@empty
	\let\@evenhead\@empty
	\def\@oddfoot{\reset@font\hfil\thepage\hfil}
	\let\@evenfoot\@oddfoot
}
\makeatother


\usepackage{amsfonts,color,morefloats,pslatex}
\usepackage{amssymb,amsthm, amsmath,latexsym}

\usepackage{algorithm}
\usepackage{algorithmic}
\usepackage{latexsym,amsmath,amsthm,amssymb,amsxtra,mathrsfs,times,CJK}
\usepackage{color}
\usepackage{amsfonts,color}
\usepackage{amssymb,amsmath,latexsym}
\usepackage{amssymb}


\usepackage[para]{threeparttable}
\def\bibsection{\section*{References}}
\usepackage{xcolor}%
\usepackage{textcomp}%
\usepackage{manyfoot}%
\usepackage{booktabs}%
\usepackage{tabularx}


\usepackage{listings}%
\usepackage{enumitem}%
\usepackage{booktabs} % for \toprule, \midrule, and \bottomrule  
\usepackage{makecell} % 用于\tabincell命令
\usepackage{caption}  % for better caption handling  
\usepackage{array}    % for \extrarowheight (not used in this example, but often useful)

\newtheorem{theorem}{Theorem}[section]
\newtheorem{lemma}[theorem]{Lemma}
\newtheorem{corollary}[theorem]{Corollary}
\newtheorem{proposition}[theorem]{Proposition}
\newtheorem{conjecture}[theorem]{Conjecture}
\newtheorem{open problem}[theorem]{Open Problem}

\newtheorem{example}[theorem]{Example}
\newtheorem{remark}[theorem]{Remark}

\newtheorem{definition}[theorem]{Definition}
\newtheorem{case}{Case}

\newcommand\myatop[2]{\genfrac{}{}{0pt}{}{#1}{#2}}
\newcommand\ov{\overline}

\newcommand{\sdef}{{\mathrm{def}}}
\newcommand{\ord}{{\mathrm{ord}}}
\newcommand{\image}{{\mathrm{Im}}}
\newcommand{\rank}{{\mathrm{rank}}}
%\newcommand{\exp}{{\mathrm{exp}}}
\newcommand{\lcm}{{\mathrm{lcm}}}
\newcommand{\tr}{{\mathrm{Tr}}}
\newcommand{\Tr}{{\mathrm{Tr}}}
\newcommand{\Norm}{{\mathrm{N}}}
\newcommand{\gf}{{\mathrm{GF}}}
\newcommand{\PG}{{\mathrm{PG}}}
\newcommand{\AG}{{\mathrm{AG}}}
\newcommand{\DAG}{{\mathrm{DAG}}}
\newcommand{\AGL}{{\mathrm{AGL}}}
\newcommand{\support}{{\mathrm{suppt}}}
\newcommand{\Aut}{{\mathrm{Aut}}}
\newcommand{\PAut}{{\mathrm{PAut}}}
\newcommand{\MAut}{{\mathrm{MAut}}}
\newcommand{\GAut}{{\mathrm{Aut}}}
\newcommand{\MAutpr}{{\mathrm{MAutpr}}}
\newcommand{\GAutpr}{{\mathrm{Autpr}}}
\newcommand{\Sym}{{\mathrm{Sym}}}
\newcommand{\rC}{{\mathrm{c}}}
\newcommand{\RS}{{\mathrm{RS}}}
\newcommand{\QRC}{{\mathrm{QRC}}}
\newcommand{\Orb}{{\mathrm{Orb}}}
\newcommand{\Stab}{{\mathrm{Stab}}}
\newcommand{\Fix}{{\mathrm{Fix}}}
\newcommand{\GL}{{\mathrm{GL}}}
\newcommand{\diff}{{\mathrm{diff}}}
\newcommand{\aut}{{\mathrm{Aut}}}
\newcommand{\rms}{{\mathrm{rms}}}

\newcommand{\wt}{{\mathtt{wt}}}
\newcommand{\hd}{{\mathtt{dist}}}
\newcommand{\RM}{{\mathrm{RM}}}

\newcommand{\F}{{\mathbb{F}}}
\newcommand{\Z}{\mathbb{{Z}}}
\newcommand{\ls}{{\mathbb{L}}}
\newcommand{\bR}{{\mathbb{R}}}
\newcommand{\bC}{{\mathbb{C}}}
\newcommand{\ef}{{\mathbb{F}}}
\newcommand{\m}{\mathbb{M}}

\newcommand{\cP}{{\mathcal{P}}}
\newcommand{\cB}{{\mathcal{B}}}
\newcommand{\cL}{{\mathcal{L}}}
\newcommand{\cI}{{\mathcal{I}}}
\newcommand{\C}{{\mathsf{C}}}
\newcommand{\M}{{\mathsf{M}}}
\newcommand{\calC}{{\mathsf{C}}}
\newcommand{\Sim}{{\mathsf{Sim}}}
\newcommand{\Ham}{{\mathsf{Ham}}}
\newcommand{\cH}{{\mathcal{H}}}
\newcommand{\cR}{{\mathcal{R}}}

\newcommand{\ba}{{\mathbf{a}}}
\newcommand{\bb}{{\mathbf{b}}}
\newcommand{\bc}{{\mathbf{c}}}
\newcommand{\bg}{{\mathbf{g}}}
\newcommand{\bm}{{\mathbf{m}}}
\newcommand{\bu}{{\mathbf{u}}}
\newcommand{\bw}{{\mathbf{w}}}
\newcommand{\bx}{{\mathbf{x}}}
\newcommand{\by}{{\mathbf{y}}}
\newcommand{\bz}{{\mathbf{z}}}
\newcommand{\tH}{{\mathbf{H}}}
\newcommand{\bI}{{\mathbf{I}}}

\newcommand{\bzero}{{\mathbf{0}}}
\newcommand{\bone}{{\mathbf{1}}}
\newcommand{\AC}{{\mathtt{AC}}}
\newcommand{\CC}{{\mathtt{CC}}}

\newcommand{\bD}{{\mathbb{D}}}

\newcommand{\dy}{{\mathrm{DY}}}
\newcommand{\rt}{{\mathrm{RT}}}
\newcommand{\p}{{\mathrm{Paley}}}

%\newcommand{\cR}{{\mathcal{R}}}
\newcommand{\cJ}{{\mathtt{J}}}
\newcommand{\cN}{{\mathcal{N}}}
%\newcommand{\AGL}{{\mathrm{AGL}}}
\newcommand{\cPGRM}{{\mathrm{PGRM}}}
\newcommand{\cGRM}{{\Omega}}


\newcommand{\ASL}{{\mathrm{ASL}}}
\newcommand{\AsiGL}{{\mathrm{A \Sigma L}}}
\newcommand{\GA}{{\mathrm{GA}}}
\newcommand{\SA}{{\mathrm{SA}}}
\newcommand{\GaA}{{\mathrm{\Gamma A}}}
\newcommand{\GaL}{{\mathrm{\Gamma L}}}
\newcommand{\SL}{{\mathrm{SL}}}
\newcommand{\PSL}{{\mathrm{PSL}}}
\newcommand{\PGL}{{\mathrm{PGL}}}
\newcommand{\PsiGL}{{\mathrm{P \Sigma L}}}
\newcommand{\PGaL}{{\mathrm{P \Gamma L}}}

\usepackage{blindtext}

\begin{document}
	
	
	\begin{frontmatter}
			\title{Beyond Polynomials: Optimal Locally Recoverable Codes from Good Rational Functions}
		
		
		\author[SWJTU]{Hengfeng Liu}
		\ead{hengfengliu@163.com}
		\author[Paris8,LAGA]{Sihem Mesnager}
		\ead{smesnager@univ-paris8.fr}
		\author[SWJTUT]{Chunming Tang\corref{cor1}}
		\ead{tangchunmingmath@163.com}
		\author[CWNU]{Xuemin Zheng}
		\ead{2455891543@qq.com}
		
		
		\cortext[cor1]{Corresponding author}
		\address[SWJTU]{School of Mathematics, Southwest Jiaotong University, Chengdu, 611756, China}
		\address[Paris8]{Department of Mathematics, University of Paris VIII, 93526 Saint-Denis, France}
		
		\address[LAGA]{CNRS, UMR 7539 LAGA, University of Paris XIII, 93430 Villetaneuse, France}
		\address[SWJTUT]{School of Information Science and Technology, Southwest Jiaotong University, Chengdu, 611756, China}
		\address[CWNU]{School of Mathematical Sciences, China West Normal University, Nanchong, 637000, China}
		
		%% Title, authors and addresses
		
		%% use the tnoteref command within \title for footnotes;
		%% use the tnotetext command for the associated footnote;
		%% use the fnref command within \author or \address for footnotes;
		%% use the fntext command for the associated footnote;
		%% use the corref command within \author for corresponding author footnotes;
		%% use the cortext command for the associated footnote;
		%% use the ead command for the email address,
		%% and the form \ead[url] for the home page:
		%%
		%% \title{Title\tnoteref{label1}}
		%% \tnotetext[label1]{}
		%% \author{Name\corref{cor1}\fnref{label2}}
		%% \ead{email address}
		%% \ead[url]{home page}
		%% \fntext[label2]{}
		%% \cortext[cor1]{}
		%% \address{Address\fnref{label3}}
		%% \fntext[label3]{}
		
		%\title{The linear codes of $2$-designs held in a class of ternary linear codes
			
		
			
		
		
	
		
		
	
		
		
		
		%% use optional labels to link authors explicitly to addresses:
		%% \author[label1,label2]{<author name>}
		%% \address[label1]{<address>}
		%% \address[label2]{<address>}
		%\author{Cunsheng Ding}
		%\ead{cding@ust.hk}
		
		
		%\cortext[lcj]{Corresponding author}
		%\address{Department of Computer Science and Engineering,
			%The Hong Kong University of Science and Technology,
			%Clear Water Bay, Kowloon, Hong Kong, China}
		
		%\tableofcontents

		
		\begin{abstract}
Locally recoverable codes (LRCs) have emerged as fundamental objects in modern coding theory, primarily due to their pivotal role in distributed and cloud storage systems. A major breakthrough in their construction was achieved by Tamo and Barg, who introduced the notion of \emph{good polynomials} as a key structural ingredient.

In this article, we propose a natural  generalization of this paradigm by introducing the concept of \emph{good rational functions}. Building upon this extension, we develop a unified and flexible framework for constructing optimal LRCs.  To quantify the quality of a rational function, we embed the problem into the rich context of algebraic function field theory and Galois theory. This perspective allows us to extend the Galois-theoretic framework originally developed by Micheli for good polynomials. In particular, we derive structural and quantitative results on the number of totally split rational places associated with rational functions. Furthermore, we construct explicit families of good rational functions that outperform all good polynomials of the same degree. As a consequence, we obtain infinite families of optimal LRCs with improved parameters compared to those arising from the classical Tamo-Barg construction. These results highlight the intrinsic strength of our approach.

\medskip

	 
	 
			
			
		\end{abstract}
		
		\begin{keyword}
			Locally recoverable code \sep good polynomial \sep rational function  \sep function field \\
			
				2020 MSC:  94B05 \sep 12E20 \sep 11C08 
		\end{keyword}
		
	\end{frontmatter}
	
	
	
	\section{Introduction}
	Let $q$ be a prime power, and let $\mathbb{F}_q$ denote the finite field with $q$ elements. Let $n,k,r$ be positive integers. An $(n,k,r)$-locally recoverable code (LRC) $\mathcal{C}$ is defined as a $k$-dimensional linear subspace of $\mathbb{F}_{q}^{n}$ with the property that any erased coordinate of a codeword $c \in \mathcal{C}$ can be recovered by accessing at most $r$ other coordinates. 
The parameter $r$ is referred to as the \emph{locality} of $\mathcal{C}$, and it measures the efficiency of the local repair process. In particular, smaller values of $r$ correspond to more efficient recovery procedures, which is a key requirement in applications such as distributed storage systems. Locally recoverable codes (LRCs) have attracted significant attention due to their fundamental role in distributed and cloud storage systems, where efficient data recovery is crucial \cite{BargTamoVladut2015}, \cite{Gopalan H S Y2012}, \cite{Micheli 2020}, \cite{Tamo Barg 2014}. Over the past decade, LRCs have been extensively studied from both theoretical and practical perspectives, leading to a rich literature on their constructions, bounds, and structural properties \cite{Cadambe Mazumdar2015}, \cite{Cai Miao Schwartz Tang}, \cite{Duke F M2022}, \cite{Jin Ma Xing 2020}, \cite{Li Ma Xing2019}, \cite{Luo Cao 2021}, \cite{Micheli 2020}, \cite{Tamo Papailiopoulos2016}, \cite{Xing Yuan 2022}, as well as recent advances \cite{Fang Tao Fu Chen Xia}, \cite{Heng2024}, \cite{Heng2025}, \cite{Luo Ezerman Ling2023}, \cite{Luo Chen Ezerman Ling2025}, \cite{Ma Xing 2023}, \cite{Sharma2025}, \cite{Tan2023}.

A fundamental limitation on the parameters of LRCs was established by Gopalan \textit{et al.} \cite{Gopalan H S Y2012}, who proved that the minimum distance $d\left (  \mathcal{C} \right )$ satisfies
\[
d\left (  \mathcal{C} \right ) \leqslant n - k - \left\lceil \frac{k}{r} \right\rceil + 2.
\]
This inequality is known as the \emph{Singleton-type bound} for LRCs. Codes attaining this bound are called \emph{optimal}, and constructing such codes with flexible parameters remains a central challenge in coding theory.

\medskip

Among the various construction paradigms, a particularly influential framework was introduced by Tamo and Barg \cite{Tamo Barg 2014}, where the notion of a \emph{good polynomial} plays a pivotal role. Intuitively, a polynomial is considered good if it exhibits a strong regularity property on suitably chosen subsets of $\mathbb{F}_q$. More precisely, a polynomial $f\in \mathbb{F}_q[X]$ is said to be $(r,l)$-good ($l>0$) if 
\begin{itemize}
	\item $f$ has degree $r+1$,
	\item there exist $l$ disjoint subsets $A_{1},\ldots, A_{l}$ in $\mathbb{F}_q$, with $|A_{i}|=r+1$ ($1\le i\le l$), such that $f$ is constant on each $A_{i}$.
\end{itemize} 

This structural property allows one to partition evaluation points into local groups, which is the key mechanism underlying locality.

\medskip

The Tamo--Barg construction is realized via generalized Reed--Solomon codes. For a message $a\in \mathbb{F}_{q}^{k}$ (written as $a_{ij}$, where $0\le i\le r-1$, $0\le j\le \frac{k}{r}-1$), the encoding polynomial is given by
\[
f_{a}(x) = \sum_{i=0}^{r-1} \sum_{j=0}^{\frac{k}{r}-1} a_{ij} f(x)^{j}x^{i}.
\]
The LRC is obtained by evaluating $f_{a}$ on the set \( A = \bigcup_{i=1}^{\ell} A_{i} \), namely,
\[
\mathcal{C} = \left\{ \big( f_{a}(x) : x \in A \big) \mid a \in \mathbb{F}_{q}^{k} \right\}.
\]

Such an \( (n, k, r) \)-LRC is optimal, with parameters $n=(r+1)l$ and $k=tr$ ($t\le l$). In addition, \cite{Tamo Barg 2014} provided explicit constructions of good polynomials arising from additive and multiplicative subgroups of $\mathbb{F}_{q}$, under certain arithmetic constraints on $(q,r)$. In particular, $r+1$ must divide either $q-1$ or $q$, or be expressible as a product of divisors of these quantities.

\medskip

These constraints motivated further generalizations. In this direction, Liu, Mesnager, and Chen \cite{Liu Mesnager Chen 2018} proposed a broader construction method based on function composition. Moreover, Liu, Mesnager, and Tang \cite{Liu Mesnager Tang 2020} showed that good polynomials can be constructed using Dickson polynomials and their compositions, thereby enriching the available families.

\medskip

To construct LRCs with a wide range of lengths and dimensions, it is crucial to identify good polynomials that are constant on as many $(r+1)$-subsets as possible. To quantify this property, for a polynomial $f\in\mathbb{F}_{q}[X]$, one defines
\begin{equation}
l_{f}=\label{size}
\#\left \{ t_{0}\in\mathbb{F}_{q}\mid f-t_{0} \; \text{has $\deg(f)$ distinct roots in $\mathbb{F}_{q}$} \right \}.
\end{equation}
This quantity measures the number of fibers of maximal size under the map induced by $f$. It is immediate that $l_{f}\le q/\deg(f)$.

\medskip

A key breakthrough in understanding this parameter was achieved by Micheli \cite{Micheli 2020}, who interpreted $l_{f}$ within the framework of algebraic function field theory. More precisely, $l_{f}$ coincides with the number $\#T^1_{split}(f)$ of totally split rational places in the global function field extension $\mathbb{F}_{q}(x)/ \mathbb{F}_{q}(t)$ defined by $f(x)=t$. Using a density argument, \cite{Micheli 2020} derived the asymptotic estimate
\[
\#T^1_{split}(f) = \frac{q}{\#G_{f}} + O(\sqrt{q}),
\]
where $G_{f}$ denotes the arithmetic monodromy group of $f$, that is, the Galois group of $f-t$ over $\mathbb{F}_{q}(t)$.

This viewpoint reveals that the construction of good polynomials is closely related to controlling the size and structure of their associated Galois groups.

\medskip

In Micheli's work \cite{Micheli 2020}, several classes of good polynomials of low degree were constructed, together with lower bounds on $l_{f}$. This line of research was further developed by Chen, Mesnager, and Zhao \cite{Chen Mesnager zhao2021}, who completely characterized good polynomials of low degree over finite fields and determined their corresponding values of $l_{f}$. Subsequently, Dukes, Ferraguti, and Micheli \cite{Duke F M2022} extended these results using an independent Galois-theoretic approach, determining $l_{f}$ for all good polynomials of degree up to $5$.

Building on these developments, Chen and Mesnager \cite{Chen Mesnager 2022} refined the Galois-theoretic framework by characterizing polynomials with minimal Galois groups and deriving explicit formulas for $l_{f}$ in several important cases, most notably for Dickson polynomials.

\medskip

These results naturally raise the question of whether broader classes of functions, beyond polynomials, can provide improved constructions of LRCs with better parameters. This question serves as one of the main motivations for the present work.


\subsection{Motivations and objectives of this paper}

Motivated by the limitations inherent in polynomial-based constructions, we propose in this paper a natural and powerful extension of the notion of good polynomial to that of \emph{good rational function}. Unlike polynomials, rational functions naturally take values on the projective line $\mathbb{P}^{1}(\mathbb{F}_{q})=\mathbb{F}_{q}\cup \left \{ \infty  \right \}$, thereby offering a richer algebraic and geometric framework.

Building upon this generalized notion, we develop a novel construction of optimal LRCs based on good rational functions. This construction preserves the fundamental principles of the Tamo--Barg framework while significantly enlarging the class of admissible functions, leading to greater flexibility in code parameters.

\medskip

A central aspect of our approach is the quantitative evaluation of how ``good " a rational function $h$ is. To this end, we investigate the number $\#T^1_{split}(h)$ of totally split rational places in the global function field extension $\mathbb{F}_{q}(x)/ \mathbb{F}_{q}(h(x))$. This perspective allows us to embed the problem into the well-established theory of algebraic function fields, thereby providing both conceptual clarity and powerful analytical tools.

In particular, our results extend and generalize those obtained by Micheli \cite{Micheli 2020} for polynomials. This extension is nontrivial, as rational functions exhibit a more intricate ramification and splitting behavior, which can be advantageously exploited.

\medskip

Moreover, in the important case where the extension $\mathbb{F}_{q}(x)/ \mathbb{F}_{q}(h(x))$ is Galois, we provide explicit characterizations of $\#T^1_{split}(h)$ from the viewpoint of group actions. This group-theoretic interpretation yields precise structural insights and enables the derivation of sharp bounds.

\medskip

In addition, we construct explicit families of good rational functions for which the associated extension $\mathbb{F}_{q}(x)/ \mathbb{F}_{q}(h(x))$ is Galois. These constructions lead to infinite families of LRCs of length $q+1$, which are particularly attractive in practice.

A striking feature of these rational functions is that they possess more totally split rational places than all good polynomials of the same degree. Consequently, the resulting LRCs achieve strictly larger lengths than those obtained via the classical Tamo--Barg construction. In a nutshell, the strength of our method stems from the following fundamental phenomenon: under suitable conditions, rational functions inherently allow for a larger number of totally split rational places than polynomials of the same degree, thereby yielding improved code constructions.

\subsection*{Organization of the paper}


The paper is organized as follows. Section \ref{pre} provides the necessary background on algebraic function fields and their extensions, which forms the theoretical foundation of our work. In particular, we review basic notions such as places, ramification, and the Hurwitz genus formula. Section \ref{LRCs from good rational functions} introduces the notion of good rational functions (Definition \ref{good rational}) and presents a general construction scheme for optimal LRCs (Theorem \ref{main construction}) based on this concept. Section \ref{how good} is devoted to the quantitative analysis of good rational functions. We relate the parameter $l$ to the number $\#T^1_{\text{split}}(h)$ of totally split rational places and study it within the framework of algebraic function field theory. Key results include Proposition \ref{split com}, Theorem \ref{split}, and Corollary \ref{estimate}. Section \ref{Gal case} focuses on the important case of Galois extensions. We develop a group-theoretic approach to characterize $\#T^1_{\text{split}}(h)$ (Theorem \ref{3 short}, Corollary \ref{good 3}) and provide explicit constructions of good rational functions via automorphism groups (Theorem \ref{gen con} and Propositions \ref{gal 1} and \ref{gal 2}).
Section \ref{comparison} presents a detailed comparison between our construction and the classical Tamo--Barg method, highlighting both the theoretical and practical advantages of our approach.
Finally, Section \ref{summary} concludes the paper with a summary of our contributions and discusses several directions for future research.



\section{Preliminaries}\label{pre}

In this section, we recall some fundamental notions and results from the theory of algebraic function fields and their extensions, which will serve as essential tools throughout the paper. We also fix notation and briefly discuss the key concepts that will be used in the subsequent sections.

\medskip

Let \( K \) be a perfect field. A rational function field over \( K \) is \( K(x) \), where \( x \) is transcendental over \( K \). A function field $F/K$ is defined as a finite algebraic extension of $K(x)$.

The set 
\[
\tilde{K}:=\{z\in F \mid z \text{ is algebraic over } K\}
\]
forms a subfield of $F$, called the \emph{constant field} of $F/K$. Clearly, one has $K \subseteq \tilde{K} \subsetneqq F$, and the extension $F/\tilde{K}$ is itself a function field over $\tilde{K}$. We say that $K$ is the \emph{full constant field} of $F$ if $K = \tilde{K}$. This condition will be particularly relevant in our later considerations involving global function fields.

\medskip

The most basic objects in the theory of function fields are valuation rings and places, which play a role analogous to primes in number fields.

A \emph{valuation ring} of \( F/K \) is a ring \( \mathcal{O} \subseteq F \) such that \( K \subsetneq \mathcal{O} \subsetneq F \) and for every \( z \in F \), either \( z \in \mathcal{O} \) or \( z^{-1} \in \mathcal{O} \). A \emph{place} \( P \) of \( F/K \) is the maximal ideal of some valuation ring \( \mathcal{O}_P \) of \( F/K \). We denote by $v_{P}$ the normalized discrete valuation associated with $P$.

Let $\mathbb{P}_F $ denote the set of all places of $F$. The quotient $O_{P}/P$ is called the \emph{residue class field} of $P$. The residue class map $O_{P}\to O_{P}/P$ induces a canonical embedding of $K$ into $O_{P}/P$, and the extension degree $[O_{P}/P:K]$ is called the \emph{degree} of $P$. A place of degree one is called a \emph{rational place} of $F/K$. These rational places will play a central role in our study of splitting behavior.

\medskip

For a rational function field \( K(x) \), the finite places correspond bijectively to irreducible monic polynomials in $K[x]$. In particular, for $a\in K$, we denote by $P_{a}$ the rational place corresponding to $x-a$. The infinite place \( P_\infty \) corresponds to the valuation associated with $1/x$.

In this paper, we focus on the case $K=\mathbb{F}_{q}$. In this setting, there are exactly $q+1$ rational places in the rational function field $\mathbb{F}_{q}(x)$. Any finite extension of $\mathbb{F}_{q}(x)$ is called a \emph{global function field}, which provides a natural analogue of number fields over finite fields.

\medskip

Let $F'/F$ be a finite separable extension of function fields. One of the fundamental phenomena in this context is the splitting of places: a place of \( F \) may decompose into several places in \( F' \), each carrying specific ramification and residue properties.

We say that a place \( P' \in \mathbb{P}_{F'} \) lies over a place \( P \in \mathbb{P}_F \), or that $P'$ is an extension of $P$, if \( P \subseteq P' \), and we write $P'|P$. For any place \( P' \in \mathbb{P}_{F'} \), there exists a unique place \( P \in \mathbb{P}_F \) such that \( P' \mid P \). Conversely, every place \( P \in \mathbb{P}_F \) has at least one, but only finitely many, places \( P' \in \mathbb{P}_{F'} \) lying over it.

\medskip

The splitting behavior of a place is described by two fundamental invariants: the ramification index and the relative degree, defined as follows.

\begin{definition}
	Let \( P' \mid P \). The ramification index \( e(P' \mid P) \) is the integer such that \( v_{P'}(z) = e(P' \mid P) \cdot v_P(z) \) for all \( z \in F \). The relative degree \( f(P' \mid P) \) is \( [\mathcal{O}_{P'}/P' : \mathcal{O}_P/P] \). We say that \( P' \mid P \) is ramified if $e(P' \mid P)>1$, and unramified if $e(P' \mid P)=1$. We say that $P$ is ramified in $F'/F$ if there exists at least one place $P'\in \mathbb{P}_{F'}$ such that $P'|P$ and $e(P' \mid P)>1$.
\end{definition}

These invariants encode the arithmetic complexity of the extension and will be crucial in analyzing the behavior of rational places in our setting.

\medskip

\begin{theorem}[Fundamental Equality]
	Let $F'/F$ be a finite separable extension of function fields and let \( P \in \mathbb{P}_F \). If \( P_1, \ldots, P_m \) are all places of \( F' \) lying over \( P \), then
	\[
	\sum_{i=1}^m e(P_i \mid P) f(P_i \mid P) = [F' : F].
	\]
\end{theorem}

This equality is a cornerstone of the theory, expressing a precise balance between ramification and residue degrees in field extensions.

\medskip

We say that a place \( P \) \emph{totally splits} in \( F'/F \) if there are exactly \( [F' : F] \) places lying over \( P \); in this case, by the fundamental equality, each of these places satisfies \( e = 1 \) and \( f = 1 \). 

An extension $P'|P$ is said to be \emph{tamely (resp., wildly) ramified} if $e(P'|P)>1$ and $\text{char}(K)$ does not divide $e(P'|P)$ (resp., $\text{char}(K)$ divides $e(P'|P)$). The extension $F'/F$ is called \emph{tame} if no place of $F$ is wildly ramified in $F'/F$.

\medskip

The Hurwitz genus formula is a powerful tool for studying the global structure of function fields and their extensions, in particular for controlling the genus and ramification.

\begin{theorem}[Hurwitz genus formula]
	Let $F'/F$ be a finite separable extension of function fields having the same constant field $K$. Let $g$ (resp. $g'$) denote the genus of $F/K$ (resp. $F'/K$). Then
	\[
	2g'-2 \geq [F' : F] \cdot (2g-2) + \sum_{P \in \mathbb{P}_F} \sum_{P' \mid P} (e(P' \mid P)-1) \cdot \deg P'.
	\]
	Equality holds if and only if $F'/F$ is tame.
\end{theorem}

This formula will play a central role in deriving bounds and structural properties in the later sections, especially in the analysis of rational functions via their associated function field extensions.


\section{LRCs from good rational functions}\label{LRCs from good rational functions}

In this section, we introduce the notion of good rational functions and show how they naturally lead to constructions of optimal locally recoverable codes. This framework can be seen as a conceptual and structural extension of the classical Tamo--Barg construction.

\medskip

Let $h(x)=\displaystyle\frac{f(x)}{g(x)}\in \mathbb{F}_q(x)$ be a rational function, with $\gcd(f(x),g(x))=1$, and degree $\deg(h)=\max\{\deg(f), \deg(g)\}$. Consider the projective line $\mathbb{P}^{1}(\mathbb{F}_{q})=\mathbb{F}_{q}\cup \left \{ \infty  \right \}$. Then $h(x)$ naturally induces a map
\[
h : \mathbb{P}^1(\mathbb{F}_q) \to \mathbb{P}^1(\mathbb{F}_q),
\]
defined by evaluation.

For any point $P \in \mathbb{P}^1(\mathbb{F}_q)$, we define the fiber
\[
h^{-1}(P) \triangleq \{ Q \in \mathbb{P}^1(\mathbb{F}_q) \mid h(Q) = P \}.
\]
By basic properties of rational functions, one always has
\[
\# h^{-1}(P) \leq \deg(h).
\]
This inequality reflects the fact that $h$ behaves as a finite morphism of degree $\deg(h)$.

\medskip

As a natural generalization of \emph{good polynomials}, we introduce the following notion of \emph{good rational function}, which will serve as the key ingredient in our construction.

\begin{definition}\label{good rational}
Let $h(x)\in \mathbb{F}_q(x)$ be a rational function, and let $l$ be a positive integer. Then $h(x)$ is said to be $(r,l)$-good if
\begin{itemize}
\item[(i)] $h(x)$ has degree $r+1$.
\item[(ii)] there exist $l$ subsets $A_{1},\ldots, A_{l}$ of $\mathbb{P}^{1}(\mathbb{F}_{q})$, such that 
\begin{itemize}
	\item $|A_{i}|=r+1$, and $A_{i}\cap A_{j}=\emptyset$, if $i\ne j$.
	\item $h(A_{i})=\{t_{i}\}$ for some $t_{i}\in\mathbb{P}^{1}(\mathbb{F}_{q}) $.
\end{itemize}
\end{itemize}
\end{definition}

This definition captures the essential feature required for locality: each subset $A_i$ behaves as a recovery set on which the function $h$ is constant.

\medskip

We say that a rational function $h$ is \emph{good} if there exists a point $P\in \mathbb{P}^{1}(\mathbb{F}_{q})$ such that $\# h^{-1}(P)= \deg(h)$; otherwise, it is said to be \emph{not good}. This condition ensures the existence of fibers of maximal size, which is crucial for constructing large families of recovery sets.

\medskip

Inspired by the Tamo--Barg framework, we now present a construction of LRCs based on $r$-good rational functions. The key idea is to exploit the structure of the fibers of $h$ to define local recovery groups.

\medskip

For a good rational function $h$, by the pigeonhole principle, there exists $b\in \mathbb{F}_{q}$ such that either $b\notin h(\mathbb{P}^{1}(\mathbb{F}_{q}))$ or $b\notin \frac{1}{h}(\mathbb{P}^{1}(\mathbb{F}_{q}))$. Since replacing $h$ by $1/h$ does not affect our construction, we may assume without loss of generality that such an element $b$ satisfies $b\notin h(\mathbb{P}^{1}(\mathbb{F}_{q}))$.

\medskip

\begin{theorem}\label{main construction}
For a $(r,l)$-good rational function $h(x)=\displaystyle\frac{f(x)}{g(x)}\in \mathbb{F}_q(x)$, pick $b\in \mathbb{F}_{q}$ such that $b\notin h(\mathbb{P}^{1}(\mathbb{F}_{q}))$. For the message $a \in \mathbb{F}_{q}^{k}$ (written as $a_{ij}$, where $0 \le i \le r-1$, $0 \le j \le \frac{k}{r}-1$), let 
\[
h_{a}(x) = \sum_{i=0}^{r-1} \sum_{j=0}^{\frac{k}{r}-1} a_{ij} \left( \frac{g(x)}{f(x)-b g(x)} \right)^{j} x^{i}.
\]

Suppose that $h=t_{i}$ on $A_{i}$ with $|A_{i}|=r+1$ ($1\le i\le l$), and let \( A = \bigcup_{i=1}^{\ell} A_{i} \).

\begin{itemize}
	\item[(i)] If $\infty \notin A$, then construct
	\[
	\mathcal{C} = \left\{ \big( h_{a}(x) : x \in A \big) \mid a \in \mathbb{F}_{q}^{k} \right\}.
	\]
	
	\item[(ii)] If $\infty \in A_{e}$ for some $1\le e\le l$, then construct
	\[
	\mathcal{C} = \Bigg\{ \underbrace{\Big( \cdots h_{a}(x) \cdots \Big)}_{x \in A\setminus\{\infty\}}, c_{\infty} \mid a \in \mathbb{F}_{q}^{k} \Bigg\},
	\]
	where $c_{\infty} = \displaystyle\sum_{j=0}^{\frac{k}{r}-1} a_{(r-1)j} \left( \frac{1}{t_{e}-b} \right)^{j}$ is the coefficient of the term $x^{r-1}$ if $h_{a}$ is viewed as a polynomial on $A_{e}$.
\end{itemize}

Let $n=(r+1)l$. Then \( \mathcal{C} \) is an optimal \( (n, k, r) \)-LRC over \( \mathbb{F}_{q} \).
\end{theorem}

\begin{proof}
The encoding is clearly linear. We now analyze its structural properties.

\medskip

For $m=1,\cdots,l$, on $A_{m}$ we have
\[
\frac{g(x)}{f(x)-b g(x)}=\frac{1}{t_{m}-b}\triangleq d_{m},
\]
and thus
\[
h_{a}(x)= \sum_{i=0}^{r-1} \sum_{j=0}^{\frac{k}{r}-1} a_{ij} d_{m}^j x^{i}.
\]

\medskip

We first prove that the mapping $a \mapsto  h_{a}(x)$ is injective. To this end, it suffices to show that the $k$ rational functions
\[
\left\{  \left( \frac{g(x)}{f(x)-b g(x)} \right)^{j}x^i: i=0,\ldots,r-1; j=0,\ldots,t-1 \right\}
\]
are linearly independent over $\mathbb{F}_q$.

Suppose that there exist coefficients $\lambda_{ij} \in \mathbb{F}_q$ such that
\[
\sum_{i=0}^{r-1} \sum_{j=0}^{\frac{k}{r}-1} \lambda_{ij} \left( \frac{g(x)}{f(x)-b g(x)} \right)^{j} x^{i}=0.
\]

Evaluating this identity on each $A_{m}$ ($1\le m\le l$), we obtain
\[
\sum_{i=0}^{r-1} \left(\sum_{j=0}^{\frac{k}{r}-1} \lambda_{ij} d_{m}^{j}\right) x^{i}=0.
\]

Since the left-hand side is a polynomial of degree at most $r-1$ that vanishes on $r+1$ distinct points, it must be identically zero. Hence,
\[
\sum_{j=0}^{\frac{k}{r}-1} \lambda_{ij} d_{m}^{j}=0, \quad 0\le i\le r-1,\; 1\le m\le l.
\]

For each fixed $i$, this yields a homogeneous linear system whose coefficient matrix is a Vandermonde matrix:
\[
 \begin{pmatrix}
	1 & d_1 & d_1^2 & \cdots & d_1^{k/r-1} \\
	1 & d_2 & d_2^2 & \cdots & d_2^{k/r-1} \\
	\vdots & \vdots & \vdots & \vdots & \vdots \\
	1 & d_l & d_l^2 & \cdots & d_l^{k/r-1}
\end{pmatrix}.
\]

Since the $d_m$ are pairwise distinct, this matrix has full rank, and therefore all $\lambda_{ij}=0$. This proves injectivity.

\medskip

We now estimate the degree of $h_{a}$. A direct computation shows that
\begin{equation}\label{degree}
(r+1)\left(\frac{k}{r}-1\right)+r-1=k+\frac{k}{r}-2\le n-2.
\end{equation}

Hence, $h_{a}$ has at most $n-2$ zeros, which implies that $\mathcal{C}$ is a $k$-dimensional subspace of $\mathbb{F}_{q}^{n}$, regardless of whether $\infty\in A$ or not.

Using (\ref{degree}), we obtain
\[
d\ge n-\deg(h_{a})= n-k-\frac{k}{r}+2,
\]
which meets the Singleton-type bound, thus proving optimality.

\medskip

We now establish the locality property.

If $\infty\notin A$, the result follows directly from Lagrange interpolation on each $A_{i}$.

If $\infty\in A_{e}$, assume the $r$ values of $h_{a}$ on $ A_{e}\setminus\{\infty\}$ are $c_{e_{1}},\cdots, c_{e_{r}}$. We show that each symbol in $\{c_{e_{1}},\cdots, c_{e_{r}},c_{\infty}\}$ can be recovered from $r$ others.

Let 
\[
b_{i}=\sum_{j=0}^{\frac{k}{r}-1} a_{ij} \left( \frac{1}{t_{e}-b} \right)^{j}.
\]
Then on $A_{e}$ we have
\[
h_{a}(x)=\sum_{i=0}^{r-1}b_{i}x^{i}, \quad \text{with } b_{r-1}=c_{\infty}.
\]

\begin{case}
$c_{\infty}$ is erased. The $r$ known evaluations determine the polynomial via interpolation, hence $c_{\infty}$ is recovered.
\end{case}

\begin{case}
$c_{v}$ is erased ($1\le v\le r$). Subtracting the leading term yields a polynomial of degree at most $r-2$, which can be reconstructed from $r-1$ values, allowing recovery of $c_v$.
\end{case}

This completes the proof.
\end{proof}

	
	
	\section{Good rational functions and algebraic function field theory}\label{how good}

Under the construction framework developed in Theorem \ref{main construction}, in order to obtain LRCs with a large range of possible lengths and dimensions, it is crucial to identify good rational functions for which the parameter $l$ is as large as possible.

\medskip

By Definition \ref{good rational}, for a rational function $h(x)$, the maximal value of $l$ is precisely given by
\begin{equation}\label{inv}
\#\left\{P \in \mathbb{P}^1(\mathbb{F}_q)\mid 	\# h^{-1}(P)=\deg (h)\right\}.
\end{equation}

In other words, $l$ is determined by the number of fibers of maximal size under the map induced by $h$. This observation naturally leads to a deeper investigation of the structure of these fibers.

\medskip

To analyze the size of the set in (\ref{inv}), we embed the problem into the framework of algebraic function field theory, where powerful tools are available to study such questions. We begin by establishing the general setting.

\medskip

Let $h(x)=\displaystyle\frac{f(x)}{g(x)}\in \mathbb{F}_q(x)$ be a rational function with $\gcd(f(x),g(x))=1$, and set $h(x)=t$. Then the extension $\mathbb{F}_q(x)/\mathbb{F}_q(t)$ is a function field extension of degree $\deg(h)$.

Throughout this paper, we impose the condition 
\[
f'(X)g(X)-g'(X)f(X)\notin \mathbb{F}_{q}[X^{q}],
\]
which ensures that the polynomial $f(X)-tg(X)$ is separable over $\mathbb{F}_{q}(t)[X]$. Consequently, the extension $\mathbb{F}_q(x)/\mathbb{F}_q(t)$ is separable. This separability condition is essential for applying the standard machinery of algebraic function field theory.

\medskip

We denote by $M_{h}$ the splitting field of the polynomial $f(X)-tg(X)$, that is, the Galois closure of the extension $\mathbb{F}_q(x)/\mathbb{F}_q(t)$. This field will play a central role in our analysis, as it allows us to interpret splitting phenomena in a Galois-theoretic framework.

\medskip

We identify each point in $ \mathbb{P}^{1}(\mathbb{F}_{q})$ with a rational place of the rational function field, and denote by $P_{a}$ the place corresponding to $a\in \mathbb{P}^{1}(\mathbb{F}_{q})$.

\medskip

The following fundamental observation establishes a precise correspondence between the fibers of $h$ and the splitting behavior of places. This connection will serve as the cornerstone of our approach.

\begin{proposition}\label{split com} 
Let $h(x)$ be a rational function, and let $a,b\in \mathbb{P}^{1}(\mathbb{F}_{q})\setminus \left \{ \infty  \right \}$, with $t=h(x)$. Then $h(a)=b$ if and only if the place $P_{a}$ in $\mathbb{F}_{q}(x)$ lies over the place $P_{b}$ in $\mathbb{F}_{q}(t)$. 

In particular, $\# h^{-1}(b)=\deg(h)$ if and only if $P_{b}$ splits totally in the extension $\mathbb{F}_q(x)/\mathbb{F}_q(t)$.
\end{proposition} 

\begin{proof}
For any $b\in\mathbb{P}^{1}(\mathbb{F}_{q})$, the rational place $P_{b}$ consists of rational functions having $b$ as a zero. If $h(a)=b$, then for any function $f(t)\in P_{b}$ in $\mathbb{F}_{q}(t)$, we have
\begin{equation}
f(h(a))=f(b)=0.
\end{equation}
This implies that $f(t)=f(h(x))\in P_{a}$ in $\mathbb{F}_{q}(x)$. Hence, $P_{b}\subseteq P_{a}$, that is, $P_{a}$ lies over $P_{b}$.

Conversely, suppose that $P_{a}$ lies over $P_{b}$, and assume for contradiction that $h(a)=b'\ne b$. By the same reasoning, $P_{a}$ would lie over $P_{b'}$, contradicting the uniqueness of the place below $P_{a}$. Therefore, $h(a)=b$.

\medskip

For the second statement, observe that $\# h^{-1}(b)=\deg(h)$ if and only if there are exactly $\deg(h)$ distinct rational places lying over $P_{b}$. By definition, this is equivalent to $P_{b}$ splitting totally in the extension $\mathbb{F}_q(x)/\mathbb{F}_q(t)$.
\end{proof}
	
	
	
	By Proposition \ref{split com}, the set (\ref{inv}) can be interpreted as the set of totally split rational places in the extension $\mathbb{F}_q(x)/\mathbb{F}_q(h(x))$. We denote this set by $T^1_{split}(h)$. 

\medskip

Therefore, our main objective is to estimate the quantity $\#T^1_{split}(h)$, which plays a central role in determining how good a rational function is in the sense of Definition \ref{good rational}. 

\medskip

A key idea in our approach is to transfer the problem to a Galois setting, where the structure of the extension becomes more tractable. Indeed, by the following standard result (see \cite{Stichtenoth2009}), it suffices to study the splitting behavior in the Galois closure $M_{h}/\mathbb{F}_q(t)$.

\begin{lemma}\label{M}
Let $h(x)$ be a rational function, and let $t=h(x)$. Let $M_{h}$ be the Galois closure of the extension $\mathbb{F}_q(x)/\mathbb{F}_q(t)$. Then for any place $P$ in $\mathbb{F}_{q}(t)$, we have
\begin{itemize}
\item[(i)] $P$ splits totally in the extension $\mathbb{F}_q(x)/\mathbb{F}_q(t)$ if and only if it splits totally in the extension $M_{h}/\mathbb{F}_q(t)$. 
\item[(ii)] $P$ is ramified in the extension $\mathbb{F}_q(x)/\mathbb{F}_q(t)$ if and only if it is ramified in the extension $M_{h}/\mathbb{F}_q(t)$. 
\end{itemize}
\end{lemma}

This lemma shows that passing to the Galois closure does not alter the essential splitting and ramification properties. Consequently, the study of $\#T^1_{split}(h)$ can be carried out in the more structured and symmetric setting of a Galois extension.

\medskip

We now establish a necessary condition for a rational function to be good, expressed in terms of the constant field of the Galois closure. This condition will be assumed throughout the remainder of the paper.

\begin{proposition}\label{not good}
If a rational function $h$ is good, then the constant field of $M_{h}$ is $\mathbb{F}_{q}$.
\end{proposition}

\begin{proof}
Assume that $h$ is good. Then there exists a rational place $P\in \mathbb{P}^{1}(\mathbb{F}_{q})$ such that $\# h^{-1}(P)= \deg(h)$. 

By Proposition \ref{split com} and Lemma \ref{M}, the place $P$ splits totally in the extension $M_{h}/\mathbb{F}_{q}(t)$. 

\medskip

Since $\mathbb{F}_{q}$ is the constant field of $\mathbb{F}_{q}(t)$, it is naturally contained in $M_{h}$. Let $P'$ be any place of $M_{h}$ lying over $P$. Because $P$ splits totally, we have $f(P'|P)=1$, that is,
\begin{equation}
\left [ O_{P'}/P':O_{P}/P\right ]=1.
\end{equation}

It follows that the residue class fields satisfy
\[
O_{P'}/P'\cong O_{P}/P\cong \mathbb{F}_{q}.
\]

\medskip

On the other hand, the constant field of $M_{h}$ is canonically embedded into each residue field $O_{P'}/P'$. Therefore, it must be a subfield of $\mathbb{F}_{q}$, and hence coincides with $\mathbb{F}_{q}$.

\medskip

This completes the proof.
\end{proof}
	
	We now further analyze the structure of the Galois function field extension $M_{h}/\mathbb{F}_q(t)$, which provides the appropriate framework for studying the splitting behavior of places.

\medskip

Let $G_{h}=\mathrm{Gal}(M_{h}/\mathbb{F}_q(t))$ be the Galois group of the extension. Let $P$ be a place of $\mathbb{F}_q(t)$. For any extension of $P$ to $M_{h}$, the ramification index $e(P)$ and the relative degree $f(P)$ depend only on $P$ (and not on the chosen extension). Moreover, the Galois group $G_{h}$ acts transitively on the set of places of $M_{h}$ lying over $P$.

\medskip

Let $P'$ be a place of $M_{h}$ lying over $P$. The stabilizer
\[
D(P'|P):=\{\sigma\in G_{h}\mid \sigma(P')=P'\}
\]
is called the \emph{decomposition group} of $P'$ over $P$. Its cardinality is given by $|D(P'|P)|=e(P)f(P)$.

Similarly, the subgroup
\[
I(P'|P):=\{\sigma\in G_{h}\mid\sigma(z)-z\in P', \ \forall z\in O_{P'}\}
\]
is called the \emph{inertia group} of $P'$ over $P$, and has order $e(P)$. It is a normal subgroup of $D(P'|P)$.

\medskip

The following classical result describes the relationship between these groups and the induced Galois action on residue fields.

\begin{theorem}\cite{Stichtenoth2009}
In the extension $M_{h}/\mathbb{F}_q(t)$, let $P'$ be a place lying over a place $P$ of $\mathbb{F}_q(t)$. Then the residue class extension $O_{P'}/P':O_{P}/P$ is Galois. Each automorphism $\sigma\in D(P'|P)$ induces an automorphism $\overline{\sigma}$ in $\mathrm{Gal}(F'_{P'}/F_{P})$ by setting
\[
\overline{\sigma}(z(P'))=\sigma(z)(P') \quad \text{for } z\in O_{P'}.
\]
The mapping
\[
D(P'|P)\longrightarrow \mathrm{Gal}(F'_{P'}/F_{P}), \quad \sigma \longmapsto \overline{\sigma}
\]
is a surjective homomorphism whose kernel is the inertia group $I(P'|P)$. In particular, one has the isomorphism
\[
D(P'|P)/I(P'|P)\cong \mathrm{Gal}(F'_{P'}/F_{P}).
\]
\end{theorem}

\medskip

We now specialize to the case where $P$ is a rational place of $\mathbb{F}_{q}(t)$. In this situation, one has $F_{P}=\mathbb{F}_{q}$, and the residue field extension $F'_{P'}/\mathbb{F}_{q}$ is finite and Galois.

\medskip

For an automorphism $\sigma\in G_{h}$ and a rational place $P$ of $\mathbb{F}_{q}(t)$, we say that $\sigma\in D(P'|P)$ is a \emph{Frobenius} for $P$ if the induced automorphism $\overline{\sigma}$ acts as
\[
\overline{\sigma}(x)=x^q,
\]
that is, $\overline{\sigma}$ generates the cyclic group $\mathrm{Gal}(F'_{P'}/\mathbb{F}_{q})$.

This notion provides a natural link between the arithmetic of the base field and the Galois action on places. In particular, it is easy to verify that $f(P)=1$ if and only if the identity automorphism acts as a Frobenius for $P$.

\medskip

We are now in a position to derive a formula for $\# T^1_{split}(h)$, which is the central quantity of interest. The proof relies on analytic tools from the theory of global function fields, namely the \emph{Zeta function} and \emph{Artin $L$-series}, which encode the distribution of places in Galois extensions.

\medskip

Denote by $R_{h}$ (resp., $R_{h}^{1}$) the set of all ramified places (resp., rational ramified places) of $\mathbb{F}_q(t)$ in the extension $M_{h}/\mathbb{F}_{q}(t)$.

\begin{theorem}\label{split}
Let $h\in\mathbb{F}_q(x)$ be a rational function, and let $M_{h}$ be the Galois closure of the extension $\mathbb{F}_q(x)/ \mathbb{F}_q(t)$. Assume that the constant field of $M_{h}$ is $\mathbb{F}_q$. Denote by $G_{h}$ the Galois group of $M_{h}/ \mathbb{F}_q(t)$.

Let $P^{1}_{M}$ be the set of rational places in $M_{h}$, and let $R_{h}^{1}$ be the set of rational ramified places of $\mathbb{F}_q(t)$ in the extension $\mathbb{F}_q(x)/ \mathbb{F}_q(t)$. Then one has
\begin{equation}
\begin{aligned}
\# T^1_{split}(h)&=\frac{\#P^1_M}{\#G_h}-\sum_{\substack{P \in R_h^1\\ f(P)=1}} \frac{1}{e(P)},
\end{aligned}
\end{equation}
where $e(P)$ and $f(P)$ denote the ramification index and the relative degree of the place $P$, respectively.
\end{theorem}
	
	\begin{proof}

We analyze the Galois extension $M_{h}/\mathbb{F}_{q}(t)$ by exploiting the deep connection between its Zeta function, that of the rational function field $\mathbb{F}_q(t)$, and the Artin $L$-series associated with irreducible characters of the Galois group $G_{h}$. This connection provides a powerful tool to count rational places.

\medskip

More precisely, one has the following factorization:
\begin{equation}\label{product}
\zeta_{M_{h}}(s)=\zeta_{\mathbb{F}_q(t)}(s)\prod_{\substack{\chi\ne \chi_{0} \\ \chi\in Irr(G_{h})}}L(s,\chi)^{\chi(1)},
\end{equation}
where $\chi$ runs through all non-trivial irreducible characters of $G_h$. We refer to \cite{Rosen 2013} for further details on this decomposition.

\medskip

Taking logarithms on both sides of (\ref{product}), we obtain
\begin{equation}\label{log}
\log \zeta_{M_{h}}(s)=\log \zeta_{\mathbb{F}_q(t)}(s) +\sum_{\substack{\chi\ne \chi_{0} \\ \chi\in Irr(G_{h})}} \chi(1)\log L(s,\chi).
\end{equation}

Let $u=q^{-s}$. Using the identity $-\log(1-u)=\sum_{m=1}^{\infty}\frac{u^m}{m}$, each term in (\ref{log}) can be expanded as a power series in $u$, yielding
\begin{equation}
\sum_{m=1}^{\infty}\frac{N_{m}(M_h)}{m}u^{m}
=
\sum_{m=1}^{\infty}\frac{N_{m}(\mathbb{F}_q(t))}{m}u^{m}
+
\sum_{\substack{\chi\ne \chi_{0} \\ \chi\in Irr(G_{h})}}
\chi(1)\sum_{P\in \mathbb{P}_{\mathbb{F}_{q}(t)}}\sum_{m=1}^{\infty}\frac{\chi(P)^{m}}{m} u^{m\deg P},
\end{equation}
where $N_{m}(L)$ denotes the number of places of degree $m$ in the function field $L$.

\medskip

By comparing the coefficients of $u$ on both sides, we obtain
\begin{equation}\label{uuuu}
\begin{aligned}
\#P_M^1
&=(q+1)+\sum_{\substack{\chi\ne \chi_{0} \\ \chi\in Irr(G_{h})}}\chi(1) \sum_{P \in  \mathbb{P}^{1}(\mathbb{F}_{q}) } \chi(P)\\
&=\sum_{P \in \mathbb{P}^{1}(\mathbb{F}_{q})}\sum_{\substack{\chi\ne \chi_{0} \\ \chi\in Irr(G_{h})}}  \chi(1)\chi(P),
\end{aligned}
\end{equation}
where
\begin{equation}
\chi(P)=\left\{
\begin{array}{cl}
\chi(\tau_P), &  P \text{ is unramified}, \\
\displaystyle\frac{1}{e(P)}\sum_{\omega \in I(P'|P)} \chi(\omega \tau_P), &  P \text{ is ramified}.
\end{array} \right.
\end{equation}
Here $P'$ is an extension of $P$ to $M_{h}$ and $\tau_P$ is a Frobenius element for $P$ in $D(P'|P)$.

\medskip

Since 
\[
q+1=\sum_{P \in \mathbb{P}^{1}(\mathbb{F}_{q})}\chi_0(1)\chi_0(P),
\]
we can combine this with (\ref{uuuu}) to obtain
\begin{equation}
\# P_M^1= \sum_{P \in  \mathbb{P}^{1}(\mathbb{F}_{q}) } \sum_{\chi\in Irr(G_{h})}\chi(1)\chi(P).
\end{equation}

\medskip

We now analyze separately the contributions of unramified and ramified places.

\medskip

\noindent
\textbf{Unramified places.}
For an unramified place $P \in  \mathbb{P}^{1}(\mathbb{F}_{q})$, we use the character identity associated with the regular representation:
\begin{equation}
\sum_{\chi\in Irr(G_{h}) }\chi(1)\chi(P)
=
\sum_{\chi\in Irr(G_{h}) }\chi(1)\chi(\tau_P)
=
\chi_{reg}(\tau_P)
=
\left\{
\begin{array}{cl}
0, &  \tau_P \neq 1,  \\
\#G_{h}, & \tau_P = 1.
\end{array}
\right.
\end{equation}

Recall that $P$ splits totally if and only if it is unramified and $\tau_P = 1$. Hence, the contribution of totally split places yields
\begin{equation}\label{half}
\begin{aligned}
\#P_M^1
=
\#T^1_{split}(h)\cdot \#G_{h}
+
\sum_{P\in R_{h}^{1}} \sum_{\chi\in Irr(G_{h})} \chi (1) \chi(P).
\end{aligned}
\end{equation}

\medskip

\noindent
\textbf{Ramified places.}
For a ramified place $P\in R_{h}^1$, one obtains
\begin{equation}\label{hhalf}
\sum_{\chi\in Irr(G_{h}) }\chi(1)\chi(P)
=
\frac{1}{e(P)}\sum_{\omega \in I(P'|P)} \sum_{\chi\in Irr(G_{h}) }\chi (1) \chi(\omega \tau_P)
=
\left\{
\begin{array}{cl}
0, &  f(P) \neq 1,  \\
\displaystyle\frac{\#G_{h}}{e(P)}, & f(P) = 1.
\end{array}
\right.
\end{equation}

\medskip

Finally, combining (\ref{half}) and (\ref{hhalf}), we obtain
\begin{equation}\label{fi}
\begin{aligned}
\# T^1_{split}(h)
=
\frac{\#P^1_M}{\#G_h}
-
\sum_{\substack{P \in R_h^1\\ f(P)=1}} \frac{1}{e(P)}.
\end{aligned}
\end{equation}

This completes the proof.
\end{proof}


We now combine Theorem \ref{split} with the classical Hasse--Weil bound to derive explicit estimates for $\#T^1_{split}(h)$. This provides a quantitative understanding of how the arithmetic of the function field $M_h$ influences the quality of the rational function $h$.

\medskip

We first recall the well-known Hasse--Weil bound (see \cite{Stichtenoth2009}).

\begin{lemma}{(Hasse-Weil bound)}\label{hasse}
Let $F/\mathbb{F}_{q}$ be a function field of genus $g$, and let $N_{1}(F)$ be the number of rational places. Then 
\begin{equation}
|N_{1}(F)-(q+1)| \leq 2g\sqrt{q}.
\end{equation}  
\end{lemma}

\medskip

Applying this bound to the function field $M_h$, and combining it with Theorem \ref{split}, we immediately obtain the following estimate.

\begin{corollary}\label{estimate}
Let $M_{h}$ be the Galois closure of $\mathbb{F}_q(x)/\mathbb{F}_q(t)$ with Galois group $G_{h}$. If the constant field of $M_{h}$ is $\mathbb{F}_q$, then
\begin{equation}\label{es}
\begin{gathered}
\left \lceil   \frac{q+1-2 g_h \sqrt{q}}{\# G_{h}}-\frac{\# R^1_h }{2}\right \rceil
\leq \# T_{\text {split }}^1(h) \leq \left \lfloor   \frac{q+1+2 g_h \sqrt{q}}{\# G_{h}} \right \rfloor,
\end{gathered}
\end{equation}
where $g_h$ is the genus of $M_{h}$, and $R^1_h$ is the set of all rational ramified places of $\mathbb{F}_q(t)$ in the extension $M_{h}/ \mathbb{F}_q(t)$.
\end{corollary}

\begin{proof}
By the Hasse--Weil bound (Lemma \ref{hasse}), we have
\[
q+1-2g_h\sqrt{q} \le \#P_{M}^{1} \le q+1+2g_h\sqrt{q}.
\]

On the other hand, for any ramified place $P\in R_{h}^{1}$, one has $e(P)\ge 2$, and hence
\[
\frac{1}{e(P)} \le \frac{1}{2}.
\]

Substituting these bounds into Equation (\ref{fi}) yields the desired inequalities in (\ref{es}).
\end{proof}

In Corollary \ref{estimate}, the lower bound of $\# T_{\text {split }}^1(h)$ depends on the parameters $g_{h}$, $\# R^1_h$, and $\# G_{h}$. Therefore, in the selection of good rational functions, it is desirable to identify those for which this lower bound is as large as possible.

\medskip

To achieve this, we derive upper bounds for $g_{h}$ and $\# R^1_h$ by employing the Riemann--Hurwitz formula (see \cite{Stichtenoth2009}). A key step in this analysis is to understand the splitting behavior of the place at infinity $P_{\infty}$ in the extension $\mathbb{F}_{q}(x)/ \mathbb{F}_{q}(t)$. Although this can be obtained directly, we state it explicitly for completeness.

\begin{proposition}\label{infi}
Let $h(x)=\displaystyle\frac{f(x)}{g(x)}$ be a rational function with $f(x)$ and $g(x)$ coprime. Let the factorization of $g(x)$ in $\mathbb{F}_{q}\left [ X \right ] $ be $g(x)=\displaystyle\prod_{ i }g_{{i}}(x)^{a_{i}}$, where the $g_{i}(x)$ are irreducible. Denote by $P_{g_{i}}$ the place associated with $g_{i}(x)$ in $\mathbb{F}_{q}(x)$, and let $P_{\infty}^{t}$ (resp., $P_{\infty}^{x}$ ) be the place at infinity in $\mathbb{F}_{q}(t)$ (resp., $\mathbb{F}_{q}(x)$). 

Let $h(x)=t$. In the extension $\mathbb{F}_{q}(x)/\mathbb{F}_{q}(t)$, the places of $\mathbb{F}_{q}(x)$ lying over $P_{\infty}^{t}$, together with their ramification indices and relative degrees, are given as follows:

\begin{itemize}
\item[(i)] If $\deg(f)\le \deg(g)$, then the places lying over $P_{\infty}^{t}$ are the $P_{g_{i}}$, with $e(P_{g_{i}}|P_{\infty}^{t})=a_{i}$ and $f(P_{g_{i}}|P_{\infty}^{t})=\deg(g_{i})$.
\item[(ii)] If $\deg(f)>\deg(g)$, then the places lying over $P_{\infty}^{t}$ are the $P_{g_{i}}$ and $P_{\infty}^{x}$, with $e(P_{g_{i}}|P_{\infty}^{t})=a_{i}$, $f(P_{g_{i}}|P_{\infty}^{t})=\deg(g_{i})$, $e(P_{\infty}^{x}|P_{\infty}^{t})=\deg(f)-\deg(g)$, and $f(P_{\infty}^{x}|P_{\infty}^{t})=1$. 
\end{itemize} 
\end{proposition}

\medskip

We are now in a position to derive explicit bounds on the number of ramified places and the genus of the Galois closure.

\begin{theorem} \label{21e}
Let $h(x)=\displaystyle\frac{f(x)}{g(x)}$ be a rational function with $f(x)$ and $g(x)$ coprime. Let $t=h(x)$, and let $M_{h}$ be the Galois closure of the extension $\mathbb{F}_q(x)/\mathbb{F}_q(t)$. In the extension $M_{h}/\mathbb{F}_q(t)$, let $R_{h}$ (resp., $R_{h}^{1}$) be the set of all ramified places (resp., ramified rational places). Let $g_{h}$ be the genus of $M_{h}$, and let $G_{h}$ be the Galois group of $M_{h}/\mathbb{F}_q(t)$. Then 

\begin{itemize}
\item[(i)] $\#R_{h}^{1}\le \#R_{h}\le \deg(h)+\displaystyle\sum_{i}\deg(g_{i})+\delta_{g}^{f}-1$, where $\deg(h)=\max\{deg(f), deg(g)  \}$, and
$$\delta_{g}^{f}=\begin{cases}
1 & \deg(f)>\deg(g), \\
0 & \deg(f)\le\deg(g).
\end{cases}$$

\item[(ii)] Suppose $\text{char}(\mathbb{F}_{q})\nmid \#G_{h}$ and the constant field of $M_{h}$ is $\mathbb{F}_{q}$. Then
$$g_{h}\le (\deg(h)-2)\#G_{h}+1.$$
\end{itemize}
\end{theorem}

\begin{proof}
We begin with the proof of $(i)$.

\medskip

By Lemma \ref{M}, the number $\#R_{h}$ coincides with the number of ramified places in the extension $\mathbb{F}_q(x)/\mathbb{F}_q(t)$. Let $\mathbb{P}(\mathbb{F}_{q})$ denote the set of all places of $\mathbb{F}_{q}(t)$. Applying the Riemann--Hurwitz formula to the extension $\mathbb{F}_q(x)/\mathbb{F}_q(t)$, we obtain
\begin{equation}\label{H1}
\begin{aligned}
2\deg(h)-2
\ge
\sum_{P\in \mathbb{P}(\mathbb{F}_{q}) }\sum_{P'|P}(e(P'|P)-1)\deg(P')
=
\sum_{P\in \mathbb{P}(\mathbb{F}_{q}) }\delta_{P}.
\end{aligned}
\end{equation}

Separating the contribution of the infinite place, we write
\[
\sum_{P\in \mathbb{P}(\mathbb{F}_{q}) }\delta_{P}
=
\sum_{\substack{P\ne P_{\infty} \\ \text{$P$ ramified}}} \delta_{P}
+
\delta_{\infty}.
\]

Using Proposition \ref{infi}, we compute
\begin{equation}\label{H2}
\begin{aligned}
\delta_{\infty}
&=
\sum_{P'|P_{\infty}}(e(P'|P_{\infty})-1)f(P'|P_{\infty}) \\
&=
\sum_{P'|P_{\infty}}e(P'|P_{\infty})f(P'|P_{\infty})
-
\sum_{P'|P_{\infty}}f(P'|P_{\infty}) \\
&=
\deg(h)-\sum_{i}\deg(g_{i})-\delta_{g}^{f}.
\end{aligned}
\end{equation}

Combining (\ref{H1}) and (\ref{H2}), we obtain
\[
\#R_{h}
\le
\sum_{\substack{P\ne P_{\infty} \\ \text{$P$ ramified}}} \delta_{P}+1
\le
2\deg(h)-\delta_{\infty}-1
=
\deg(h)+\sum_{i}\deg(g_{i})+\delta_{g}^{f}-1.
\]

This proves $(i)$.

\medskip

We now prove $(ii)$.

\medskip

We apply the Riemann--Hurwitz formula to the Galois extension $M_{h}/\mathbb{F}_q(t)$. Since $\text{char}(\mathbb{F}_{q})\nmid \#G_{h}$ and the constant field of $M_{h}$ is $\mathbb{F}_{q}$, the extension is tame, and equality holds. Hence,
\begin{equation}\label{ii1}
\begin{aligned}
2g_{h}-2
&=
-2\#G_{h}
+
\sum_{P\in \mathbb{P}(\mathbb{F}_{q}) }\deg(P)\sum_{P'|P}(e(P'|P)-1)f(P'|P).
\end{aligned}
\end{equation}

Using the transitivity of the Galois action, we obtain
\[
\sum_{P'|P}(e(P'|P)-1)f(P'|P)
=
(e(P)-1)f(P)\cdot \frac{\#G_{h}}{e(P)f(P)}.
\]

Substituting into (\ref{ii1}), we deduce
\[
2g_{h}-2
\le
\#G_{h}\left ( \sum_{P\in \mathbb{P}(\mathbb{F}_{q}) }\deg(P)-2 \right ).
\]

To bound $\sum_{P}\deg(P)$, we use again (\ref{H1}), which yields
\begin{equation}\label{ii2}
\sum_{P\in \mathbb{P}(\mathbb{F}_{q}) }\deg(P)
\le
\sum_{P\in \mathbb{P}(\mathbb{F}_{q}) }\delta_{P}
\le
2\deg(h)-2.
\end{equation}

Combining (\ref{ii1}) and (\ref{ii2}), we conclude that
\[
g_{h}\le (\deg(h)-2)\#G_{h}+1.
\]

This completes the proof.
\end{proof}


Using the results in Corollary \ref{estimate} and Theorem \ref{21e}, we can derive an explicit lower bound on the maximal value of $l$ for a good rational function. This provides a concrete criterion for selecting functions that yield LRCs with large parameters.

\begin{proposition}\label{pregene}
Let $h(x)=\displaystyle\frac{f(x)}{g(x)}$ be a rational function with $f(x)$ and $g(x)$ coprime, and let $\deg(h)=r+1$. Let $t=h(x)$, and let $M_{h}$ be the Galois closure of the extension $\mathbb{F}_q(x)/\mathbb{F}_q(t)$. Suppose that the constant field of $M_{h}$ is $\mathbb{F}_{q}$ and $\gcd\left ( q,(r+1)! \right )=1 $. If $h$ is $(r,l)$-good, then $l$ satisfies
\[
l \ge \frac{q-2\sqrt{q}+1}{(r+1)!}-2\sqrt{q}(r-1)-\frac{r+1+\sum_{i}\deg(g_{i})}{2},
\]
where $g_{i}$ runs over all irreducible factors of $g(x)$ over $\mathbb{F}_{q}\left [ X \right ]$.
\end{proposition}

\begin{proof}
Since the polynomial $f(X)-tg(X)$ has degree $r+1$, its Galois group $G_{h}$ is a subgroup of the symmetric group $S_{r+1}$, and hence $\#G_{h}\mid (r+1)!$.

The condition $\gcd\left ( q,(r+1)! \right )=1$ implies that $\text{char}(\mathbb{F}_{q})\nmid \#G_{h}$. Therefore, the assumptions of Theorem \ref{21e} are satisfied.

Applying Theorem \ref{21e} together with Corollary \ref{estimate}, we immediately obtain the stated lower bound on $l$.
\end{proof}

\medskip

We now discuss the notion of equivalence between rational functions, which is important for identifying genuinely new constructions.

\medskip

For two rational functions $f,g\in \mathbb{F}_{q}(x)$, it is well known that $\mathbb{F}_{q}(f(x))=\mathbb{F}_{q}(g(x))$ if and only if $f$ is obtained from $g$ by left composition with a linear fractional transformation. More precisely, $f=\phi \circ g$, where $\phi=\displaystyle \frac{ax+b}{cx+d}$ with $ad-bc\ne 0$.

In this case, we say that $f$ and $g$ are \emph{equivalent}, since they define the same function field extension $\mathbb{F}_{q}(x)/\mathbb{F}_{q}(t)$. Consequently, in the search for good rational functions, it suffices to consider representatives up to this equivalence relation.

\medskip

In particular, we are primarily interested in rational functions that are \emph{not equivalent to polynomials}, as these yield genuinely new constructions beyond the classical Tamo--Barg framework.

\medskip

The following general construction provides explicit examples of such rational functions.

\begin{proposition}
Let $r$ be a positive integer satisfying $\gcd\left ( q,(r+1)! \right )=1$, and let $S\subseteq  \mathbb{F}_{q}$ with $\#S=r+1$. Let 
\[
h(x)=\displaystyle \frac{\prod_{s\in S}\left ( x-s \right )}{x-a},
\]
where $a\notin S$. Then $h$ is $(r,l)$-good with
\[
l\ge \frac{q-2\sqrt{q}+1}{(r+1)!}-2\sqrt{q}(r-1)-\frac{r+2}{2}.
\]
\end{proposition}

\begin{proof}
The rational function $h$ is clearly good, since $\#h^{-1}(0)=r+1$. 

By Proposition \ref{not good}, the constant field of $M_{h}$ is $\mathbb{F}_{q}$. Therefore, all the conditions of Proposition \ref{pregene} are satisfied, and the desired lower bound follows directly.
\end{proof}

\medskip

\begin{remark}
The function $h=\displaystyle \frac{\prod_{s\in S}\left ( x-s \right )}{x-a}$ is not equivalent to any polynomial under linear fractional transformations.

Indeed, $h$ has two distinct poles, namely $\{a, \infty\}$. If $h=\phi \circ g$, where $\phi(x)=\displaystyle \frac{ax+b}{cx+d}$ and $g$ is a polynomial, then $g(\infty)=\infty$, and $g(a)$ would correspond to two distinct poles of $\phi$, which is impossible.

This shows that the above construction genuinely produces rational functions outside the polynomial framework.
\end{remark}


\section{The case of Galois extension}\label{Gal case}

In our code construction, a $(r,l)$-good rational function $h$ produces locally recoverable codes (LRCs) of length $(r+1)l$, where $l=\# T_{\text{split}}^1(h)$. This parameter $l$ plays a central role in determining the length and efficiency of the resulting codes. 

It is worth emphasizing that Corollary \ref{estimate} has already shown that the quantity $\# T_{\text{split}}^1(h)$ is approximately of order $q/\#G_{h}$. This asymptotic behavior naturally suggests that, for a fixed locality parameter $r$, one should seek rational functions of degree $r+1$ whose associated Galois group $G_h$ is as small as possible. Indeed, minimizing the size of $G_h$ leads to a maximization of the number of totally split places, and therefore to longer codes.

Recall that $G_{h}$ is a transitive subgroup of the symmetric group $S_{r+1}$. In particular, the extremal case occurs when $G_{h}$ reaches its minimal possible size, namely $\#G_h = r+1$, which happens if and only if the extension $\mathbb{F}_{q}(x)/ \mathbb{F}_{q}(t)$ is Galois, where $h(x)=t$. This characterization provides a natural and structurally rich setting in which the interplay between algebraic function fields and group actions becomes particularly apparent.

In view of these observations, it is both natural and fruitful to focus on the case where the extension $\mathbb{F}_{q}(x)/\mathbb{F}_{q}(t)$ is Galois. This setting allows us to exploit the full symmetry encoded by the Galois group and to obtain a more precise understanding of the splitting behavior of rational places.

When the extension $\mathbb{F}_{q}(x)/\mathbb{F}_{q}(t)$ is Galois, the group $G_{h}$ admits a natural action on the set of rational places of $\mathbb{F}_{q}(x)$. This group action provides a powerful tool to analyze the arithmetic and geometric properties of the function $h$. In particular, the orbit structure induced by $G_h$ encodes the splitting behavior of places in the extension, and hence directly influences the cardinality of $T_{\text{split}}^1(h)$.

In this context, we propose an explicit characterization of $\# T_{\text {split }}^1(h)$ from the perspective of group actions. This approach not only clarifies the combinatorial nature of the splitting phenomenon but also reveals that the upper bound for $\# T_{\text{split}}^1(h)$ given in Corollary \ref{estimate} is in fact sharp in the Galois case.

Beyond this structural analysis, we also develop general and explicit constructions of good rational functions arising from subgroups of the automorphism group of $\mathbb{F}_{q}(x)$. These constructions naturally yield Galois extensions of the form $\mathbb{F}_{q}(x)/\mathbb{F}_{q}(t)$ and provide a systematic framework for generating suitable functions $h$.

As a consequence, infinite families of LRCs of length $q+1$ are obtained directly from these good rational functions via Theorem \ref{main construction}. These families exhibit particularly attractive parameters and structural properties. The advantages of these constructions, especially in comparison with the classical Tamo--Barg LRCs, will be highlighted and analyzed in detail in the next section. In particular, we will see how the Galois-theoretic perspective leads to improvements in both flexibility and performance.



\subsection{Characterization of Galois extension}

We now exploit the Galois structure of the extension to obtain a precise description of the splitting behavior via group actions.

\medskip

Consider the natural action of $G_{h}$ on the set of rational places of $\mathbb{F}_{q}(x)$. We classify the orbits according to their size: an orbit is said to be \emph{long} if its size is $\#G_{h}=\deg(h)$, and \emph{short} otherwise.

\medskip

The following result shows that short orbits are extremely limited in number.

\begin{theorem}\label{3 short}
Let $h(x)$ be a rational function such that the extension $\mathbb{F}_{q}(x)/\mathbb{F}_{q}(t)$ is Galois, where $t=h(x)$. Let $G_{h}$ be the Galois group. Consider its natural action on the rational places of $\mathbb{F}_{q}(x)$. Then 
\begin{itemize}
\item[(i)] There are at most three short orbits.
\item[(ii)] The total number of orbits is $\left \lceil \frac{q+1}{\deg(h)}  \right \rceil$ when there are $0$ or $1$ short orbits, and $\left \lceil \frac{q+1}{\deg(h)}  \right \rceil+1$ when there are $2$ or $3$ short orbits.
\end{itemize}
\end{theorem}

\begin{proof}
The action of $G_{h}$ partitions the $q+1$ rational places of $\mathbb{P}^{1}(\mathbb{F}_{q})$ into $w$ disjoint orbits of sizes $l_{1}\le \cdots \le l_{w}$.

\medskip

Two rational places lie in the same orbit if and only if they lie over the same rational place of $\mathbb{F}_{q}(t)$. Hence, each orbit corresponds to a rational place $P$ of $\mathbb{F}_{q}(t)$, and its size is given by
\[
l_i = \frac{\deg(h)}{e(P)},
\]
where $e(P)$ is the ramification index.

\medskip

Let the first $k$ orbits be short. Then $e(P)>1$ for these orbits. Applying the Riemann--Hurwitz formula, we obtain
\begin{equation}\label{2d-2}
\begin{aligned}
2\deg(h)-2
&\ge
\sum_{P\in \mathbb{P}^{1}(\mathbb{F}_{q}) }\sum_{P'|P}(e(P)-1)\deg(P') \\
&\ge
\sum_{\substack{P \text{ rational ramified}}}
\sum_{P'|P}(e(P)-1) \\
&=
\sum_{i=1}^{k}\frac{\deg(h)}{e(P)}(e(P)-1) \\
&=
\sum_{i=1}^{k}\left ( \deg(h)-l_{i}\right ) \\
&\ge
k\cdot\frac{\deg(h)}{2}.
\end{aligned}
\end{equation}

This implies $k\le 3$, proving $(i)$.

\medskip

We now turn to $(ii)$. Clearly, the number of orbits satisfies
\[
w \ge \left \lceil \frac{q+1}{\deg(h)} \right \rceil.
\]

We distinguish cases.

\begin{itemize}
\item If $w=1$, then $q+1=l_{1}\le \deg(h)$, hence $\left \lceil \frac{q+1}{\deg(h)} \right \rceil=1$, and there are at most one short orbit.

\item If $w=2$, then $q+1=l_{1}+l_{2}$, and one checks that $w=2=\left \lceil \frac{q+1}{\deg(h)} \right \rceil+1$ if and only if both orbits are short.
\end{itemize}

Now assume $w\ge 3$. Let the first $k\le 3$ orbits be short. From (\ref{2d-2}), we deduce that
\[
l_{1}+l_{2}+l_{3}\ge \deg(h)+2.
\]

Therefore,
\begin{equation}\label{os}
\begin{aligned}
\frac{q+1}{\deg(h)}
=
\frac{l_{1}+l_{2}+l_{3}+\deg(h)(w-3)}{\deg(h)}
\ge
w-2+\frac{2}{\deg(h)}.
\end{aligned}
\end{equation}

This yields
\[
\left \lceil \frac{q+1}{\deg(h)} \right \rceil \le w \le \left \lceil \frac{q+1}{\deg(h)} \right \rceil +1.
\]

Moreover, equality $w=\left \lceil \frac{q+1}{\deg(h)} \right \rceil+1$ holds if and only if
\[
l_{1}+l_{2}+l_{3}\le 2\deg(h),
\]
which is equivalent to having $2$ or $3$ short orbits.

This completes the proof.
\end{proof}

\medskip

We now deduce an explicit formula for $\# T_{\text{split}}^1(h)$.

\begin{corollary}\label{good 3}
Let $h(x)$ be a rational function such that the extension $\mathbb{F}_{q}(x)/\mathbb{F}_{q}(t)$ is Galois. Let $V_{s}$ denote the number of short orbits under the action of $G_{h}$. Then
\[
\# T_{\text{split}}^1(h) =
\begin{cases}
\dfrac{q+1}{r+1}, & \text{if } V_s = 0, \\[2ex]
\left\lceil \dfrac{q+1}{r+1} \right\rceil - 1, & \text{if } V_s = 1,2, \\[2ex]
\left\lceil \dfrac{q+1}{r+1} \right\rceil - 2, & \text{if } V_s = 3.
\end{cases}
\]
\end{corollary}

\begin{proof}
By definition, $\# T_{\text{split}}^1(h)$ is equal to the number of long orbits. The result follows directly from Theorem \ref{3 short}$(ii)$.
\end{proof}

\medskip

\begin{remark}\label{remark bound}
The upper bound of $\# T_{\text {split }}^1(h)$ given in Corollary \ref{estimate} is 
\[
\left \lfloor   \frac{q+1}{r+1} \right \rfloor.
\]
By Corollary \ref{good 3}, if $r+1\nmid q+1$, then
\[
\left \lfloor \frac{q+1}{r+1} \right \rfloor
=
\left \lceil \frac{q+1}{r+1} \right \rceil -1.
\]
Hence, the upper bound is achieved whenever there are one or two short orbits.
\end{remark}

\subsection{Constructions of good rational functions}

We now present a general framework for constructing good rational functions that induce Galois extensions $\mathbb{F}_{q}(x)/\mathbb{F}_{q}(h(x))$. These constructions yield infinite families of LRCs whose lengths exceed those obtained from polynomials of the same degree.

\medskip

Our approach relies on the rich group-theoretic structure of the automorphism group of the rational function field $\mathbb{F}_{q}(x)$. It is well known that this group is given by
\[
\left\{\frac{ax+b}{cx+d}\mid ad-bc\ne 0\right\},
\]
which is isomorphic to the projective general linear group $\mathrm{PGL}(2, q)$ via the correspondence
\[
\frac{ax+b}{cx+d}\mapsto
\begin{pmatrix}
a & b\\
c & d
\end{pmatrix}.
\]
In what follows, we identify these two structures.

\medskip

The group $\mathrm{PGL}(2, q)$ acts on $\mathbb{F}_{q}(x)$ by right composition: for $h\in \mathbb{F}_{q}(x)$ and $\phi\in \mathrm{PGL}(2, q)$, we define
\[
\phi \cdot h = h \circ \phi.
\]

\medskip

Let
\[
Z = \left\{ f \in \mathbb{F}_{q}(x) : f \circ \phi = f \text{ for all } \phi \in\mathrm{PGL}(2, q) \right\}
\]
be the fixed field of $\mathrm{PGL}(2, q)$. Then $\mathbb{F}_{q}(x)/Z$ is a Galois extension with Galois group $\mathrm{PGL}(2, q)$.

\medskip

More generally, for any subgroup $H \le \mathrm{PGL}(2, q)$, we denote by $\mathbb{F}_q(x)^{H}$ its fixed field. Then $\mathbb{F}_q(x)/\mathbb{F}_q(x)^{H}$ is a Galois extension, and by the Galois correspondence one has
\[
\mathrm{Gal}(\mathbb{F}_q(x)/\mathbb{F}_q(x)^{H})=H.
\]

\medskip

By L\"uroth's theorem (see \cite{Stichtenoth2009}), any intermediate field of $\mathbb{F}_q(x)$ is rational. Hence there exists a function $h_H(x)\in \mathbb{F}_q(x)$ such that
\[
\mathbb{F}_q(x)^{H}=\mathbb{F}_q(h_H(x)).
\]
Moreover, $h_H(x)$ is uniquely determined by $H$ up to left composition by an element of $\mathrm{PGL}(2, q)$, and its degree satisfies
\[
\deg (h_H(x))=[ \mathbb{F}_q(x):\mathbb{F}_q(h_H(x))]=\#H.
\]

\medskip

This observation provides a direct and systematic way to construct rational functions with prescribed Galois group. In particular, for locality parameter $r$, if there exists a subgroup $H$ of $\mathrm{PGL}(2, q)$ of order $r+1$, then there exists a rational function $h_H(x)$ of degree $r+1$ such that the extension $\mathbb{F}_{q}(x)/\mathbb{F}_{q}(t)$ (with $t=h_H(x)$) is Galois.

\medskip

The group $\mathrm{PGL}(2, q)$ has order $q(q^2-1)$ and possesses a rich subgroup structure (see \cite{Dickson1958, Leemans 2009}). In particular, when $r+1$ divides one of $q-1$, $q$, or $q+1$, the group $\mathrm{PGL}(2, q)$ contains a subgroup of order $r+1$.

\medskip

Moreover, when $r+1$ divides $q-1$ or $q+1$, or when $r+1=p$ (where $q=p^s$ and $p$ is prime), $\mathrm{PGL}(2, q)$ contains a cyclic subgroup of order $r+1$. This case is especially important, as cyclic subgroups lead to particularly explicit constructions.

\medskip

Based on this observation, we first present in Theorem \ref{gen con} a general construction of good rational functions arising from cyclic subgroups of $\mathrm{PGL}(2, q)$. 

\medskip

We then specialize to low-degree cases, namely degrees $3$ and $4$, for which we provide explicit formulas and compare their performance with that of polynomials of the same degree.

\medskip

We do not consider degree $2$, since quadratic polynomials used in the Tamo--Barg construction already attain the upper bound in Corollary \ref{estimate}.


\begin{theorem}\setcounter{case}{0} \label{gen con}
Let $\phi \in \mathrm{PGL}(2,q)\setminus \mathrm{AGL}(1,q)$ be a non-affine transformation ($\phi(x) = \displaystyle\frac{ax+b}{cx+d}$ with $c \neq 0$, $ad-bc\ne 0$), and suppose that the order of $\phi$ is $r+1$. Let $\phi^{i}=\underbrace{\phi \circ\cdots \circ \phi}_{i}$ (with $\phi^{0}=x$), and define
\[
h(x)=\sum_{i=0}^{r}\phi^{i}(x).
\]
Then $\deg(h)=r+1$, and the fixed field of the cyclic subgroup $\langle \phi \rangle$ is precisely $\mathbb{F}_{q}(h(x))$. In particular, the extension $\mathbb{F}_{q}(x)/\mathbb{F}_{q}(h(x))$ is Galois.
\end{theorem}

\begin{proof}

We begin by showing that for every integer $1 \leq i \leq r$, the transformation $\phi^i$ is not affine. Assume, for contradiction, that $\phi^i$ is affine for some $i$. Then $\phi^i$ fixes the point at infinity. 

If $\phi^i$ is a translation, that is, $\phi^i(x)=x+b$ with $b\neq 0$, then $\infty$ is its unique fixed point. Since $\phi$ commutes with $\phi^i$, it must preserve this fixed point, and hence $\phi(\infty)=\infty$, which contradicts the assumption that $\phi$ is non-affine.

If $\phi^i(x)=ax+b$ with $a\neq 1$, then $\phi^i$ has exactly two fixed points: $\infty$ and a finite point $P$. Since $\phi$ commutes with $\phi^i$, it permutes the set $\{\infty,P\}$. If $\phi$ fixes both points, then $\phi(\infty)=\infty$, again a contradiction. If $\phi$ swaps them, then $\phi$ must have order $2$, which contradicts the assumption that its order is $r+1$. 

Therefore, all iterates $\phi^i$ for $1 \le i \le r$ are non-affine.

\medskip

We now determine the degree of $h$. For $i \geq 1$, write
\[
\phi^i(x)=\frac{a_i x+b_i}{c_i x+d_i}, \quad \text{with } c_i \neq 0,
\]
and set $\phi^0(x)=x$. Consider the polynomial
\[
D(x)=\prod_{i=1}^r (c_i x+d_i),
\]
which has degree $r$. Then $h(x)$ can be written as
\[
h(x)=\frac{N(x)}{D(x)},
\]
where
\[
N(x)=xD(x)+\sum_{i=1}^r (a_i x+b_i)\frac{D(x)}{c_i x+d_i}.
\]

The term $xD(x)$ has degree $r+1$, while each summand in the second term has degree at most $r$. It follows that
\[
\deg(N)=r+1.
\]

\medskip

We next show that $N(x)$ and $D(x)$ are coprime. Suppose that $x_0$ is a common root. Then $D(x_0)=0$, so there exists some index $j$ such that $c_j x_0+d_j=0$. 

For any $i\neq j$, one must have $c_i x_0+d_i\neq 0$, since otherwise $x_0$ would be a pole of both $\phi^i$ and $\phi^j$. This would imply that
\[
\phi^{j-i}(\infty)=\infty,
\]
which contradicts the fact established above that no nontrivial iterate of $\phi$ is affine.

Consequently,
\[
N(x_0)=(a_j x_0+b_j)\prod_{i\ne j}(c_i x_0+d_i)\neq 0,
\]
which is a contradiction. Hence $\gcd(N(x),D(x))=1$, and therefore $\deg(h)=r+1$.

\medskip

Finally, we establish the Galois property. By construction, the function $h$ is invariant under $\phi$, and therefore under the entire cyclic group $\langle \phi \rangle$. This yields the inclusion
\[
\mathbb{F}_q(h(x)) \subseteq \mathbb{F}_q(x)^{\langle \phi \rangle}.
\]

Since
\[
[\mathbb{F}_q(x):\mathbb{F}_q(h(x))]=\deg(h)=r+1=\#\langle \phi \rangle,
\]
it follows that $\mathbb{F}_q(h(x))$ coincides with the fixed field of $\langle \phi \rangle$.

\medskip

This proves that the extension $\mathbb{F}_q(x)/\mathbb{F}_q(h(x))$ is Galois with Galois group $\langle \phi \rangle$.

\end{proof}





Based on Theorem \ref{gen con}, we now construct explicit rational functions $h$ for which the extension $\mathbb{F}_{q}(x)/\mathbb{F}_{q}(h(x))$ is Galois.

\medskip

We recall that $\mathrm{PGL}(2, q)$ acts naturally on the projective line $\mathbb{P}^{1}(\mathbb{F}_{q})$. More precisely, for $u\in \mathbb{P}^{1}(\mathbb{F}_{q})$ and $\phi=\displaystyle\frac{ax+b}{cx+d}$, the action is given by
\[
u \longmapsto \frac{au+b}{cu+d}.
\]

\medskip

We now present an explicit construction in the case of degree $3$.

\begin{proposition}\label{gal 1}
Let $q$ be a prime power, $w\in\mathbb{F}_{q}^{*}$, and define
\[
h(x)=\displaystyle\frac{x^3-3w^2x+w^3}{x(x-w)}.
\]
Then the extension $\mathbb{F}_{q}(x)/\mathbb{F}_{q}(h(x))$ is Galois and 
\begin{itemize}
\item[(i)] If $q=3^m$, then $\# T_{\text{split}}^1(h) =\displaystyle\frac{q}{3}$, and $h$ is $(2, \displaystyle\frac{q}{3})$-good.

\item[(ii)] If $q\equiv 1 \pmod{3}$, then $\# T_{\text{split}}^1(h) =\displaystyle\frac{q-1}{3}$, and $h$ is $(2, \displaystyle\frac{q-1}{3})$-good.

\item[(iii)] If $q\equiv 2 \pmod{3}$, then $\# T_{\text{split}}^1(h) =\displaystyle\frac{q+1}{3}$, and $h$ is $(2, \displaystyle\frac{q+1}{3})$-good.
\end{itemize}
\end{proposition}

\begin{proof}

By Theorem \ref{gen con}, it suffices to consider the function
\[
h(x)=\sum_{i=0}^{2}\phi^{i}(x),
\]
where $\phi$ is an element of order $3$ in $\mathrm{PGL}(2,q)$.

Let
\[
\phi(x)=\frac{w^2}{w-x}.
\]
A direct computation shows that
\[
\phi^2(x)=\frac{w(x-w)}{x}, \qquad \phi^3(x)=x.
\]
Consequently,
\[
x+\frac{w^2}{w-x}+\frac{w(x-w)}{x}
=
\frac{x^3-3w^2x+w^3}{x(x-w)}.
\]
It follows that $h(x)$ is invariant under the cyclic group $\langle \phi \rangle$, and therefore the extension $\mathbb{F}_{q}(x)/\mathbb{F}_{q}(h(x))$ is Galois with Galois group $G_h=\langle \phi \rangle$.

\medskip

We now determine the short orbits of the action of $\langle \phi \rangle$. These correspond precisely to the fixed points of $\phi$, that is, to the solutions of
\begin{equation}\label{fixpoint}
x = \frac{w^2}{w-x} \quad \Longleftrightarrow \quad x^2 - wx + w^2 = 0.
\end{equation}
Thus, the number of short orbits is governed by the number of solutions of this quadratic equation in $\mathbb{F}_q$.

\medskip

We proceed by distinguishing the different cases according to the characteristic of the field.

\smallskip

\emph{Case 1: $q=2^m$.}  
Let $\operatorname{Tr}: \mathbb{F}_{2^m}\to\mathbb{F}_2$ denote the trace map. It is well known that equation (\ref{fixpoint}) has a solution in $\mathbb{F}_{2^m}$ if and only if $\operatorname{Tr}(1)=0$, which is equivalent to $m$ being even.

If $m$ is even, then $q\equiv 1 \pmod{3}$ and there is exactly one short orbit. If $m$ is odd, then $q\equiv 2 \pmod{3}$ and no short orbit occurs. Applying Corollary \ref{good 3} yields the corresponding values of $\#T_{\text{split}}^1(h)$.

\medskip

\emph{Case 2: $q$ odd.}

In this case, the discriminant of equation (\ref{fixpoint}) is $-3w^2$. If $q=3^m$, then the discriminant vanishes, and the equation admits a unique solution. This gives rise to a single short orbit and hence $\#T_{\text{split}}^1(h)=\frac{q}{3}$.

Assume now that $3\nmid q$. Then equation (\ref{fixpoint}) has two distinct solutions if and only if $-3$ is a square in $\mathbb{F}_q$. Accordingly, there are two short orbits when $-3$ is a quadratic residue, and no short orbit otherwise.

By quadratic reciprocity, $-3$ is a square in $\mathbb{F}_q$ if and only if $q\equiv 1 \pmod{3}$, which completes the case analysis.

\medskip

In each of the above situations, applying Corollary \ref{good 3} yields the claimed formulas for $\#T_{\text{split}}^1(h)$.

\end{proof}






\begin{proposition}\label{gal 2}
Let $q$ be an odd prime power, $d\in\mathbb{F}_{q}^{*}$, and define
\[
h(x)=\displaystyle\frac{4x^4-12d^{2}x^{2}+8d^{3}x-d^{4}}{2x(x-d)(2x-d)}.
\]
Then the extension $\mathbb{F}_{q}(x)/\mathbb{F}_{q}(h(x))$ is Galois and 
\begin{itemize}
\item[(i)] If $q\equiv 3 \pmod{4}$, then $\# T_{\text{split}}^1(h) =\displaystyle\frac{q+1}{4}$, and $h$ is $(3, \displaystyle\frac{q+1}{4})$-good.

\item[(ii)] If $q\equiv 1 \pmod{4}$, then $\# T_{\text{split}}^1(h) =\displaystyle\frac{q-1}{4}$, and $h$ is $(3, \displaystyle\frac{q-1}{4})$-good.
\end{itemize}
\end{proposition}


\begin{proof}

Let
\[
\phi(x)=\frac{d^2}{2(d-x)}.
\]
A direct computation shows that
\[
\phi^2(x)=\frac{d(d-x)}{d-2x}, \quad 
\phi^3(x)=-\frac{d(d-2x)}{2x}, \quad 
\phi^4(x)=x.
\]
It follows that $\phi$ has order $4$, and one readily verifies that
\[
h(x)=\sum_{i=0}^{3}\phi^{i}(x).
\]
By Theorem \ref{gen con}, the extension $\mathbb{F}_{q}(x)/\mathbb{F}_{q}(h(x))$ is Galois with Galois group $G_h=\langle \phi \rangle$.

\medskip

We now analyze the orbit structure under the action of $G_h$. Since $\#G_h=4$, the possible orbit sizes are $1,2,$ and $4$, and the short orbits are precisely those of size $1$ or $2$.

Orbits of size $1$ correspond to fixed points of $\phi$, that is, to the solutions of
\begin{equation}\label{fixpoint1}
x = \frac{d^2}{2(d-x)} \quad \Longleftrightarrow \quad 2x^2 - 2dx + d^2 = 0.
\end{equation}

On the other hand, orbits of size $2$ would correspond to points satisfying $\phi^2(x)=x$ but $\phi(x)\ne x$. However,
\begin{equation}\label{fixpoint2}
x = \frac{d(d-x)}{d-2x} \quad \Longleftrightarrow \quad 2x^2 - 2dx + d^2 = 0,
\end{equation}
which coincides with equation (\ref{fixpoint1}). Consequently, every solution of $\phi^2(x)=x$ is already a fixed point of $\phi$, and no orbit of size $2$ occurs. It follows that all short orbits are exactly the fixed points of $\phi$.

\medskip

We now determine the number of such fixed points by studying the solutions of
\begin{equation}\label{fix2}
2x^2 - 2dx + d^2 = 0.
\end{equation}
Since $q$ is odd, the discriminant of this quadratic equation is
\[
\Delta = (-2d)^2 - 8d^2 = -4d^2.
\]
Thus, $\Delta$ is a square in $\mathbb{F}_q$ if and only if $-1$ is a square in $\mathbb{F}_q$, which holds precisely when $q\equiv 1 \pmod{4}$.

If $q\equiv 1 \pmod{4}$, equation (\ref{fix2}) admits two distinct solutions, yielding two short orbits. If $q\equiv 3 \pmod{4}$, it has no solution, and hence no short orbit arises.

\medskip

Applying Corollary \ref{good 3}, we obtain
\[
\# T_{\text{split}}^1(h) =
\begin{cases}
\left\lceil \dfrac{q+1}{4} \right\rceil -1 = \dfrac{q-1}{4}, & \text{if } q\equiv 1 \pmod{4}, \\[2ex]
\dfrac{q+1}{4}, & \text{if } q\equiv 3 \pmod{4}.
\end{cases}
\]

This completes the proof.

\end{proof}




\begin{remark}
The above constructions illustrate the strength of our approach based on Galois extensions.

\medskip

\begin{itemize}
\item The rational functions constructed in Proposition \ref{gal 1} and Proposition \ref{gal 2} achieve the upper bound in Corollary \ref{estimate}, showing that this bound is sharp.

\item When $q\equiv 2 \pmod{3}$ (resp., $q\equiv 3\pmod{4}$), the functions in Proposition \ref{gal 1} (resp., Proposition \ref{gal 2}) yield infinite families of LRCs of length $q+1$ with locality $2$ (resp., $3$).

\item In particular, these constructions provide longer codes than those obtained via the Tamo--Barg polynomial framework for the same locality.
\end{itemize}

\end{remark}

\section{Comparison with the Tamo--Barg construction}\label{comparison}

In this section, we compare the LRCs obtained from good rational functions with those arising from the classical Tamo--Barg construction based on good polynomials.

\medskip

Recall that in the Tamo--Barg framework, the length of the resulting code is determined by the number of disjoint $(r+1)$-subsets on which a good polynomial is constant. Equivalently, this corresponds to the number of totally split rational places associated with the polynomial. As shown in \cite{Micheli 2020}, this quantity is typically of order
\[
\frac{q}{\#G_f} + O(\sqrt{q}),
\]
where $G_f$ is the arithmetic monodromy group of the polynomial $f$.

\medskip

Our construction replaces polynomials with rational functions and exploits the structure of Galois extensions. As established in Corollary \ref{estimate}, the number of totally split rational places satisfies
\[
\# T_{\text{split}}^1(h) \approx \frac{q}{\#G_h},
\]
with explicit control on the error term.

\medskip

\noindent
\textbf{Fundamental difference.}

The key distinction lies in the size of the Galois group. For polynomials of degree $r+1$, the Galois group is typically large (often close to the full symmetric group $S_{r+1}$), leading to
\[
\#G_f \approx (r+1)!,
\]
and consequently a relatively small number of totally split places.

\medskip

By contrast, our construction allows one to realize rational functions with Galois group of size exactly $r+1$, which is the smallest possible for a transitive subgroup. This leads to a significantly larger number of totally split rational places.

\medskip

\noindent
\textbf{Quantitative comparison.}

For a fixed locality parameter $r$, the Tamo--Barg construction typically yields
\[
l \approx \frac{q}{(r+1)!},
\]
whereas our construction achieves
\[
l \approx \frac{q}{r+1}.
\]

\medskip

Thus, the improvement factor is roughly $(r)!$, which becomes substantial even for moderate values of $r$.

\medskip

\noindent
\textbf{Optimality.}

Moreover, in the Galois case, Corollary \ref{good 3} shows that our construction attains the upper bound
\[
\# T_{\text{split}}^1(h) \le \left\lfloor \frac{q+1}{r+1} \right\rfloor,
\]
demonstrating that our approach is optimal with respect to this parameter.

\medskip

\noindent
\textbf{Concrete gains.}

The explicit constructions given in Propositions \ref{gal 1} and \ref{gal 2} yield infinite families of optimal LRCs of length
\[
n = (r+1)\cdot \# T_{\text{split}}^1(h) = q+1,
\]
which is the maximum possible length for evaluation on $\mathbb{P}^{1}(\mathbb{F}_q)$.

\medskip

In particular:
\begin{itemize}
\item for locality $r=2$, we obtain codes of length $q+1$ when $q\equiv 2 \pmod{3}$;
\item for locality $r=3$, we obtain codes of length $q+1$ when $q\equiv 3 \pmod{4}$.
\end{itemize}

\medskip

\noindent
\textbf{Conceptual advantage.}

Beyond the quantitative improvement, our framework provides a conceptual shift: instead of constructing good polynomials with small monodromy groups (which is often difficult), we directly construct rational functions with controlled Galois group via group actions.

\medskip

This group-theoretic perspective offers a flexible and systematic approach, enabling the construction of new families of LRCs that are inaccessible via polynomial-based methods.

\medskip

\noindent
\textbf{Summary.}

In summary, compared with the Tamo--Barg construction, our method:
\begin{itemize}
\item achieves significantly larger code lengths for the same locality;
\item attains the theoretical upper bound on the number of totally split rational places;
\item provides explicit and infinite families of optimal LRCs;
\item introduces a new and versatile framework based on rational functions and Galois theory.
\end{itemize}





\section{Summary and concluding remarks}\label{summary}

In this paper, we have extended the notion of good polynomials to the broader and more flexible framework of good rational functions, leading to a new construction paradigm for optimal locally recoverable codes (LRCs). 

\medskip

Building upon the foundational Tamo--Barg construction, our approach replaces polynomials with rational functions defined on the projective line $\mathbb{P}^1(\mathbb{F}_q)$. This extension allows for a substantially richer structure, enabling the construction of codes with improved parameters. In particular, we introduced the concept of $(r,l)$-good rational functions and developed an explicit LRC construction scheme achieving the Singleton-type bound.

\medskip

A central contribution of this work is the identification of a deep connection between the quality of a rational function and the arithmetic of the associated function field extension $\mathbb{F}_q(x)/\mathbb{F}_q(h(x))$. By interpreting the parameter $l$ as the number $\# T_{\text{split}}^1(h)$ of totally split rational places, we were able to apply tools from algebraic function field theory and Galois theory to obtain both asymptotic estimates and exact formulas.

\medskip

In the important case where the extension is Galois, we derived explicit characterizations of $\# T_{\text{split}}^1(h)$ via group actions and showed that our construction achieves the theoretical upper bound. This leads to infinite families of optimal LRCs with length $q+1$, which is maximal for evaluation over $\mathbb{P}^1(\mathbb{F}_q)$.

\medskip

Moreover, through concrete constructions for small degrees (notably $r=2$ and $r=3$), we demonstrated that good rational functions can systematically outperform good polynomials of the same degree. In particular, our approach yields significantly longer codes for a fixed locality, highlighting a fundamental advantage over the classical Tamo--Barg framework.

\medskip

\noindent
\textbf{Future directions.}

This work opens several promising research directions. A natural continuation is the systematic classification of good rational functions of low degree, together with a precise determination of their associated Galois groups. This direction parallels recent advances for polynomials in \cite{Chen Mesnager zhao2021} and \cite{Duke F M2022}, and may lead to a deeper understanding of the interplay between algebraic structures and coding parameters.

\medskip

Another interesting direction is to explore more general group actions beyond cyclic subgroups of $\mathrm{PGL}(2,q)$, which may yield new families of rational functions with desirable splitting properties.

\medskip

Finally, it would be of interest to investigate potential applications of good rational functions in related areas, such as algebraic geometry codes, cryptographic constructions, or pseudorandom structures, where controlled splitting behavior plays a crucial role.

\medskip

\noindent
In conclusion, the introduction of good rational functions provides both a conceptual and practical advancement in the design of locally recoverable codes, opening new avenues for further developments in coding theory and finite field arithmetic.
 


	
\begin{thebibliography}{99}

\bibitem{BargTamoVladut2015}
A. Barg, I. Tamo, and S. Vladut,
``Locally recoverable codes on algebraic curves,''
in Proc. IEEE Int. Symp. Inf. Theory (ISIT), 2015, pp. 1252--1256.

\bibitem{Cadambe Mazumdar2015}
V. R. Cadambe and A. Mazumdar,
``Bounds on the size of locally recoverable codes,''
IEEE Trans. Inf. Theory, vol. 61, no. 11, pp. 5787--5794, 2015.

\bibitem{Cai Miao Schwartz Tang}
H. Cai, Y. Miao, M. Schwartz, and X. Tang,
``On optimal locally repairable codes with super-linear length,''
IEEE Trans. Inf. Theory, vol. 66, no. 8, pp. 4853--4868, 2020.

\bibitem{Chenbo}
B. Chen, G. Zhang, and W. Li,
``New optimal linear codes with hierarchical locality,''
IEEE Trans. Inf. Theory, vol. 69, no. 3, pp. 1544--1550, 2022.

\bibitem{Chen Mesnager zhao2021}
R. Chen, S. Mesnager, and C. Zhao,
``Good polynomials for optimal LRC of low locality,''
Des. Codes Cryptogr., vol. 89, no. 7, pp. 1639--1660, 2021.

\bibitem{Chen Mesnager 2022}
R. Chen and S. Mesnager,
``A function field approach toward good polynomials for further results on optimal LRC codes,''
Finite Fields Appl., vol. 81, p. 102028, 2022.

\bibitem{Dickson1958}
L. E. Dickson,
\emph{Linear Groups: With an Exposition of the Galois Field Theory}.
New York: Dover, 1958.

\bibitem{Duke F M2022}
A. Dukes, A. Ferraguti, and G. Micheli,
``Optimal selection for good polynomials of degree up to five,''
Des. Codes Cryptogr., vol. 90, no. 6, pp. 1427--1436, 2022.

\bibitem{Fang Tao Fu Chen Xia}
W. Fang, R. Tao, F. Fu, B. Chen, and S. Xia,
``Bounds and constructions of singleton-optimal locally repairable codes with small localities,''
IEEE Trans. Inf. Theory, vol. 70, no. 10, pp. 6842--6856, 2024.

\bibitem{Gopalan H S Y2012}
P. Gopalan, C. Huang, H. Simitci, and S. Yekhanin,
``On the locality of codeword symbols,''
IEEE Trans. Inf. Theory, vol. 58, no. 11, pp. 6925--6934, 2012.

\bibitem{Heng2024}
Z. Heng, X. Li, Y. Wu, and Q. Wang,
``Two families of linear codes with desirable properties from some functions over finite fields,''
IEEE Trans. Inf. Theory, 2024.

\bibitem{Heng2025}
Z. Heng, M. Yang, and Y. Ming,
``A family of self-orthogonal divisible codes with locality 2,''
Discrete Math., vol. 348, no. 10, p. 114529, 2025.

\bibitem{Jin Ma Xing 2020}
L. Jin, L. Ma, and C. Xing,
``Construction of optimal locally repairable codes via automorphism groups of rational function fields,''
IEEE Trans. Inf. Theory, vol. 66, no. 1, pp. 210--221, 2020.

\bibitem{Leemans 2009}
D. Leemans and E. Schulte,
``Polytopes with groups of type $PGL(2,q)$,''
Ars Math. Contemp., vol. 2, no. 2, pp. 163--171, 2009.

\bibitem{Li Heng2023}
X. Li and Z. Heng,
``Constructions of near MDS codes which are optimal locally recoverable codes,''
Finite Fields Appl., vol. 88, p. 102184, 2023.

\bibitem{Li Ma Xing2019}
X. Li, L. Ma, and C. Xing,
``Optimal locally repairable codes via elliptic curves,''
IEEE Trans. Inf. Theory, vol. 65, no. 1, pp. 108--117, 2019.

\bibitem{Liu Mesnager Chen 2018}
J. Liu, S. Mesnager, and L. Chen,
``New constructions of optimal locally recoverable codes via good polynomials,''
IEEE Trans. Inf. Theory, vol. 64, no. 2, pp. 889--899, 2018.

\bibitem{Liu Mesnager Tang 2020}
J. Liu, S. Mesnager, and D. Tang,
``Constructions of optimal locally recoverable codes via Dickson polynomials,''
Des. Codes Cryptogr., vol. 88, no. 9, pp. 1759--1780, 2020.

\bibitem{Luo Cao2020}
G. Luo and X. Cao,
``Optimal cyclic codes with hierarchical locality,''
IEEE Trans. Commun., vol. 68, no. 6, pp. 3302--3310, 2020.

\bibitem{Luo Cao 2021}
G. Luo and X. Cao,
``Constructions of optimal binary locally recoverable codes via a general construction of linear codes,''
IEEE Trans. Commun., vol. 69, no. 8, pp. 4987--4997, 2021.

\bibitem{Luo Ezerman Ling2023}
G. Luo, M. F. Ezerman, and S. Ling,
``Three new constructions of optimal locally repairable codes from matrix-product codes,''
IEEE Trans. Inf. Theory, vol. 69, no. 1, pp. 75--85, 2023.

\bibitem{Luo Chen Ezerman Ling2025}
G. Luo, B. Chen, M. F. Ezerman, and S. Ling,
``Bounds and constructions of quantum locally recoverable codes from quantum CSS codes,''
IEEE Trans. Inf. Theory, vol. 71, pp. 1794--1802, 2025.

\bibitem{Ma Xing 2023}
L. Ma and C. Xing,
``The group structures of automorphism groups of elliptic curves over finite fields and their applications to optimal locally recoverable codes,''
J. Combin. Theory Ser. A, vol. 193, p. 105686, 2023.

\bibitem{Micheli 2020}
G. Micheli,
``Constructions of locally recoverable codes which are optimal,''
IEEE Trans. Inf. Theory, vol. 66, no. 1, pp. 167--175, 2020.

\bibitem{Rosen 2013}
M. Rosen,
\emph{Number Theory in Function Fields}.
Springer, 2013.

\bibitem{Sharma2025}
S. Sharma, V. Ramkumar, and I. Tamo,
``Quantum locally recoverable codes via good polynomials,''
IEEE J. Sel. Areas Inf. Theory, 2025.

\bibitem{Stichtenoth2009}
H. Stichtenoth,
\emph{Algebraic Function Fields and Codes}.
Springer, 2009.

\bibitem{Tamo Barg 2014}
I. Tamo and A. Barg,
``A family of optimal locally recoverable codes,''
IEEE Trans. Inf. Theory, vol. 60, no. 8, pp. 4661--4676, 2014.

\bibitem{Tamo Papailiopoulos2016}
I. Tamo, D. S. Papailiopoulos, and A. G. Dimakis,
``Optimal locally repairable codes and connections to matroid theory,''
IEEE Trans. Inf. Theory, vol. 62, no. 12, pp. 6661--6671, 2016.

\bibitem{Tan2023}
P. Tan, C. Fan, C. Ding, C. Tang, and Z. Zhou,
``The minimum locality of linear codes,''
Des. Codes Cryptogr., vol. 91, no. 1, pp. 83--114, 2023.

\bibitem{Xing Yuan 2022}
C. Xing and C. Yuan,
``Construction of optimal $(r,\delta)$-locally recoverable codes and connection with graph theory,''
IEEE Trans. Inf. Theory, vol. 68, no. 7, pp. 4320--4328, 2022.

\end{thebibliography}
	
\end{document}



